\documentclass{article}
\usepackage[english]{babel}
\usepackage{geometry,amsmath,amssymb,latexsym,theorem}
\geometry{letterpaper}

%%%%%%%%%% Start TeXmacs macros
\newcommand{\assign}{:=}
\newcommand{\cdummy}{\cdot}
\newcommand{\tmaffiliation}[1]{\\ #1}
\newcommand{\tmem}[1]{{\em #1\/}}
\newcommand{\tmop}[1]{\ensuremath{\operatorname{#1}}}
\newcommand{\tmscript}[1]{\text{\scriptsize{$#1$}}}
\newcommand{\tmtextbf}[1]{\text{{\bfseries{#1}}}}
\newcommand{\tmtextup}[1]{\text{{\upshape{#1}}}}
\newenvironment{proof}{\noindent\textbf{Proof\ }}{\hspace*{\fill}$\Box$\medskip}
\newtheorem{corollary}{Corollary}
\newtheorem{definition}{Definition}
\newtheorem{lemma}{Lemma}
{\theorembodyfont{\rmfamily\small}\newtheorem{problem}{Problem}}
\newtheorem{proposition}{Proposition}
{\theorembodyfont{\rmfamily}\newtheorem{remark}{Remark}}
\newtheorem{theorem}{Theorem}
%%%%%%%%%% End TeXmacs macros

\newtheorem{algorithm}{Algorithm}
\newcommand{\Prim}{\tmop{Prim}}
\newcommand{\Sym}{\tmop{Sym}}
\newcommand{\ev}{\tmop{ev}}
\newcommand{\Aut}{\tmop{Aut}}
\newcommand{\ord}{\tmop{ord}}
\newcommand{\Res}{Res}
\newcommand{\trdeg}{\tmop{trdeg}}
\newcommand{\CA}{\mathfrak{B}}
\newcommand{\C}{\mathbb{C}}
\newcommand{\Z}{\mathbb{Z}}
\newcommand{\N}{\ensuremath{\mathbb{N}}}
\newcommand{\R}{\mathbb{R}}
\newcommand{\Q}{\ensuremath{\mathbb{Q}}}
\newcommand{\Lop}{\mathcal{L}}

\begin{document}

\title{The Coefficient Algebra of Exponentials of Rational Functions}

\author{
  Stephen Crowley
  \tmaffiliation{July 3, 2026}
}

\date{}

\maketitle

\begin{abstract}
  Let $R$ be a rational function with a pole of order $m$ at $a$ and principal
  part $\Pi (w) = \sum_{j = 1}^m c_j w^{- j}$, $w = z - a$. The Laurent
  coefficients of $e^R$ on a punctured disk about $a$ are governed by the
  coefficients $\Lambda_N$ of $\exp (\Pi)$, which satisfy the first-order
  recurrence $N \Lambda_N = \sum_{j = 1}^{\min (m, N)} jc_j \Lambda_{N - j}$.
  We take this recurrence as the primitive object, construct from it a graded
  commutative Hopf algebra $\CA_m = \C [c_1, \ldots, c_m]$ carrying the
  sequence $(\Lambda_N)$, and prove its structure theory: $\CA_m$ is
  canonically the quotient of the Hopf algebra of symmetric functions by the
  Hopf ideal $(p_{m + 1}, p_{m + 2}, \ldots)$, under which $\Lambda_N$
  corresponds to the complete homogeneous symmetric function $h_N$, $c_j$ to
  $p_j / j$, and the universal recurrence to Newton's identity. The sequence
  $(\Lambda_N)$ is a divided-power sequence; as functions of $c_1$ the $N!
  \hspace{0.17em} \Lambda_N$ form an Appell sequence; the coefficient matrix
  is an element of the Appell subgroup of the exponential Riordan group; $N!
  \hspace{0.17em} \Lambda_N$ enumerates set partitions with block sizes at
  most $m$; the graded Hopf automorphism group is the torus $(\C^{\times})^m$;
  and the Hankel determinants of $(\Lambda_N)$ are, up to explicit sign,
  rectangular Schur functions. We then prove that the analytic
  Laurent-coefficient construction is a model of the algebra, with
  quantitative growth bounds; we identify the classical specializations
  exactly (modified Bessel functions for $m = 1$, Hermite-type polynomials for
  $m = 2$, Wright generalized Bessel functions for single-term principal parts
  against a linear entire exponent) and refrain from claiming reductions that
  do not exist for $m \ge 3$; we prove holonomy of the full Laurent
  coefficient sequence at every pole, with an explicit recurrence and order
  bound derived from $y' = R' y$; and we prove an initiality theorem and a
  recognition theorem that together constitute the axiomatic characterization,
  with independence of the hypotheses established by counterexamples. The
  final section develops, from the fractional Riccati equation itself
  (M{\"u}ntz--Tau series, moment functional, the polynomial
  Chebyshev--Wheeler three-term recurrence, its diagonal Pad{\'e} numerator
  and denominator, the Jacobi-matrix spectrum, and real-ray Stahl
  convergence, all proved in full), the rough Heston characteristic function
  $\varphi_M = \exp \left( t \hspace{0.17em} A_M (x) / B_M (x) \right)$, $x =
  t^{\alpha}$, as the exponential of a diagonal Pad{\'e} approximant in the
  coefficient algebra: we give the exact local Laurent data of $\varphi_M$
  at the zeros of $B_M$, prove convergence of Pad{\'e} prices to exact
  prices under explicit hypotheses on a fixed Lewis-type contour with no
  contour deformation through approximant poles, and derive and prove exact
  European call and put price formulas directly from the cumulant
  generating function via Gil--Pelaez inversion, including a certified
  real-axis pole-freeness result for the Pad{\'e} approximant's evaluation
  contour. Statements that remain open (the domination hypothesis feeding
  price convergence, resurgent/alien-derivative structure,
  uniform-in-parameter domination, exponential-motive positioning) are listed
  as open problems, not theorems.
\end{abstract}

\section{The universal recurrence}\label{sec:recurrence}

Throughout, $m \ge 1$ is a fixed integer and $c_1, \ldots, c_m$ are
independent commuting indeterminates over $\C$. Set
\[ \CA_m \assign \C [c_1, \ldots, c_m], \]
graded by $\mathrm{wt} (c_j) = j$, so that $\CA_m = \bigoplus_{N \ge 0}
(\CA_m)_N$ with $(\CA_m)_0 = \C$. All constructions in
Sections~\ref{sec:recurrence}--\ref{sec:axioms} are exact identities in
$\CA_m$ or in the formal power series ring $\CA_m [[v]]$; no analysis enters
until Section~\ref{sec:model}.

\begin{definition}
  \label{def:E}$E (v) \assign \exp \left( \sum_{j = 1}^m c_j v^j \right) \in
  \CA_m [[v]]$, and $\Lambda_N \in \CA_m$ is defined by $E (v) = \sum_{N \ge
  0} \Lambda_N v^N$. We set $\Lambda_N \assign 0$ for $N < 0$.
\end{definition}

The exponential is well defined because $\sum_j c_j v^j$ has zero constant
term. Under the substitution $v = w^{- 1}$, $E$ becomes the formal expansion
of $\exp \left( \sum_{j = 1}^m c_j w^{- j} \right)$, the exponential of the
principal part of a rational function at a pole of order $\le m$; the analytic
counterpart is Theorem~\ref{thm:analyticmodel}.

\begin{lemma}
  [Universal recurrence]\label{lem:recurrence} $\Lambda_0 = 1$ and, for all $N
  \ge 1$,
  \begin{equation}
    \label{eq:universal} N \hspace{0.17em} \Lambda_N \hspace{0.27em} =
    \hspace{0.27em} \sum_{j = 1}^{\min (m, N)} j \hspace{0.17em} c_j 
    \hspace{0.17em} \Lambda_{N - j} .
  \end{equation}
\end{lemma}

\begin{proof}
  $E$ satisfies the linear first-order equation $E' (v) = \left( \sum_{j =
  1}^m jc_j v^{j - 1} \right) E (v)$ in $\CA_m [[v]]$, obtained by formal
  differentiation of the exponential. The coefficient of $v^{N - 1}$ on the
  left is $N \Lambda_N$; on the right it is $\sum_{j = 1}^m jc_j \Lambda_{N -
  j}$, in which terms with $j > N$ vanish by the convention $\Lambda_{< 0} =
  0$. The initial value $\Lambda_0 = E (0) = 1$.
\end{proof}

\begin{lemma}
  [Partition formula]\label{lem:partition} For all $N \ge 0$,
  \begin{equation}
    \label{eq:partition} \Lambda_N \hspace{0.27em} = \hspace{0.27em}
    \sum_{\tmscript{\begin{array}{c}
      k_1, \ldots, k_m \ge 0\\
      \sum_j jk_j = N
    \end{array}}} \hspace{0.27em} \prod_{j = 1}^m \frac{c_j^{\hspace{0.17em}
    k_j}}{k_j !} .
  \end{equation}
  In particular $\Lambda_N$ is homogeneous of weight $N$.
\end{lemma}

\begin{proof}
  $E (v) = \prod_{j = 1}^m \exp (c_j v^j) = \prod_{j = 1}^m \sum_{k_j \ge 0}
  \frac{c_j^{k_j}}{k_j !} v^{jk_j}$; the product of the $m$ series is computed
  by collecting the coefficient of $v^N$, which is \eqref{eq:partition}.
  Homogeneity is immediate since the monomial $\prod c_j^{k_j}$ has weight
  $\sum jk_j = N$.
\end{proof}

\begin{proposition}
  [Bell-polynomial form]\label{prop:bell} $N! \hspace{0.17em} \Lambda_N = B_N 
  (1! \hspace{0.17em} c_1, \hspace{0.17em} 2! \hspace{0.17em} c_2, \ldots, m!
  \hspace{0.17em} c_m, \hspace{0.17em} 0, \ldots, 0)$, where $B_N$ is the
  complete Bell polynomial, defined by $\exp \left( \sum_{j \ge 1} x_j t^j /
  j! \right) = \sum_{N \ge 0} B_N  (x_1, \ldots, x_N)  \hspace{0.17em} t^N /
  N!$ {\tmem{(Comtet {\cite[{\textsection}3.3]{Comtet}})}}.
\end{proposition}

\begin{proof}
  Substitute $x_j = j! \hspace{0.17em} c_j$ for $j \le m$ and $x_j = 0$ for $j
  > m$ in the defining generating identity of $B_N$; the left side becomes $E
  (t)$.
\end{proof}

\begin{lemma}
  [Derivations]\label{lem:derivation} For $1 \le j \le m$ and all $N$,
  \[ \frac{\partial \Lambda_N}{\partial c_j} \hspace{0.27em} = \hspace{0.27em}
     \Lambda_{N - j} . \]
\end{lemma}

\begin{proof}
  $\partial_{c_j} E (v) = v^j E (v)$; compare coefficients of $v^N$.
\end{proof}

\begin{remark}
  \label{rem:eulerequiv}Lemma~\ref{lem:recurrence} is exactly the Euler
  homogeneity relation: the grading operator $\mathcal{E}= \sum_{j = 1}^m j
  \hspace{0.17em} c_j  \hspace{0.17em} \partial_{c_j}$ satisfies $\mathcal{E}
  \Lambda_N = \sum_j jc_j \Lambda_{N - j}$ by Lemma~\ref{lem:derivation}, and
  $\mathcal{E} \Lambda_N = N \Lambda_N$ by Lemma~\ref{lem:partition}.
  Conversely, any sequence of weight-homogeneous elements satisfying
  Lemma~\ref{lem:derivation} with $\Lambda_0 = 1$ satisfies
  \eqref{eq:universal}. The universal recurrence, the lowering operators, and
  the grading are one structure viewed three ways.
\end{remark}

\begin{theorem}
  [Generation and freeness]\label{thm:generation} $\C [\Lambda_1, \ldots,
  \Lambda_m] = \CA_m$, and $\Lambda_1, \ldots, \Lambda_m$ are algebraically
  independent over $\C$. Consequently $(\Lambda_N)_{N \ge 0} \subset \C
  [\Lambda_1, \ldots, \Lambda_m]$: the entire coefficient sequence lives in
  the polynomial algebra on its first $m$ terms.
\end{theorem}

\begin{proof}
  For $1 \le N \le m$, the multi-index $k_N = 1$, $k_j = 0$ $(j \ne N)$ is the
  unique solution of $\sum jk_j = N$ using the variable $c_N$; every other
  solution has $k_N = 0$ and involves only $c_1, \ldots, c_{N - 1}$. Hence
  $\Lambda_N = c_N + q_N (c_1, \ldots, c_{N - 1})$ with $q_N$ polynomial. By
  induction, $c_N \in \C [\Lambda_1, \ldots, \Lambda_N]$, so $\C [\Lambda_1,
  \ldots, \Lambda_m] \supseteq \C [c_1, \ldots, c_m]$; the reverse inclusion
  is trivial. Since $\trdeg_{\C} \CA_m = m$ and $\Lambda_1, \ldots, \Lambda_m$
  generate it, they are algebraically independent. For $N > m$, $\Lambda_N \in
  \CA_m = \C [\Lambda_1, \ldots, \Lambda_m]$; explicitly, iterate
  \eqref{eq:universal}.
\end{proof}

\section{Hopf structure and the symmetric-function
realization}\label{sec:hopf}

\begin{definition}
  \label{def:hopf}Equip $\CA_m$ with the unique bialgebra structure for which
  each $c_j$ is primitive: $\Delta c_j = c_j \otimes 1 + 1 \otimes c_j$,
  $\varepsilon (c_j) = 0$, extended as algebra morphisms. Since $\CA_m$ is
  graded and connected ($(\CA_m)_0 = \C$), it is automatically a Hopf algebra,
  with antipode determined by $S (c_j) = - c_j$ {\tmem{(graded connected
  bialgebras possess antipodes; see
  {\cite[{\textsection}1]{MilnorMoore}}{\hspace{0.17em}})}}.
\end{definition}

\begin{theorem}
  [Divided-power property]\label{thm:dividedpower} $\Delta \Lambda_N = \sum_{i
  = 0}^N \Lambda_i \otimes \Lambda_{N - i}$ for all $N \ge 0$, $\varepsilon
  (\Lambda_N) = \delta_{N, 0}$, and $\sum_{i = 0}^N \Lambda_i  \hspace{0.17em}
  S (\Lambda_{N - i}) = \delta_{N, 0}$, with $S (\Lambda_N) = \Lambda_N  (-
  c_1, \ldots, - c_m)$.
\end{theorem}

\begin{proof}
  Apply the algebra morphism $\Delta$ coefficient-wise in $(\CA_m \otimes
  \CA_m) [[v]]$:
  \[ \sum_N \Delta \Lambda_N  \hspace{0.17em} v^N = \Delta E (v) = \exp \left(
     \sum_j (c_j \otimes 1 + 1 \otimes c_j) v^j \right) = \exp \left( \sum_j
     (c_j \otimes 1) v^j \right) \exp \left( \sum_j (1 \otimes c_j) v^j
     \right), \]
  the factorization being valid because the two exponents commute. The right
  side equals $\left( \sum_i \Lambda_i \otimes 1 \hspace{0.17em} v^i \right)
  \left( \sum_k 1 \otimes \Lambda_k  \hspace{0.17em} v^k \right)$; compare
  coefficients. The counit and antipode identities follow the same way from
  $\varepsilon E (v) = 1$ and $SE (v) = E (v)^{- 1} = \exp (- \sum c_j v^j)$.
\end{proof}

Thus $(\Lambda_N)$ is a {\tmem{divided-power sequence}} in the graded
connected Hopf algebra $\CA_m$: $\Lambda_N \in (\CA_m)_N$, $\Lambda_0 = 1$,
$\Delta \Lambda_N = \sum \Lambda_i \otimes \Lambda_{N - i}$. Equivalently, the
series $E (v)$ is group-like in the completed Hopf algebra $\CA_m [[v]]$.

\begin{theorem}
  [Primitives]\label{thm:primitives} $\Prim (\CA_m) = \bigoplus_{j = 1}^m \C
  \hspace{0.17em} c_j$; in particular the primitive space in weight $j$ is
  one-dimensional for $1 \le j \le m$ and zero for $j > m$.
\end{theorem}

\begin{proof}
  Each $c_j$ is primitive by construction. Conversely, let $f = \sum_{\alpha}
  a_{\alpha} c^{\alpha}$ be primitive, with multi-index notation $c^{\alpha} =
  \prod c_j^{\alpha_j}$. Then $\Delta (c^{\alpha}) = \prod_j (c_j \otimes 1 +
  1 \otimes c_j)^{\alpha_j} = \sum_{0 \le \beta \le \alpha}
  \binom{\alpha}{\beta} c^{\beta} \otimes c^{\alpha - \beta}$. The component
  of $\Delta f - f \otimes 1 - 1 \otimes f$ in $\CA_m^+ \otimes \CA_m^+$
  ($\CA_m^+ \assign \ker \varepsilon$) is $\sum_{\alpha} a_{\alpha}  \sum_{0 <
  \beta < \alpha} \binom{\alpha}{\beta} c^{\beta} \otimes c^{\alpha - \beta}$.
  The elements $\{c^{\beta} \otimes c^{\gamma} : \beta, \gamma \ne 0\}$ are
  linearly independent in $\CA_m \otimes \CA_m$, and for fixed $(\beta,
  \gamma)$ the total coefficient is $a_{\beta + \gamma} \binom{\beta +
  \gamma}{\beta}$. Primitivity therefore forces $a_{\alpha} = 0$ whenever
  $\alpha$ admits a decomposition $\alpha = \beta + \gamma$ with $\beta,
  \gamma \ne 0$, i.e. whenever $| \alpha | \assign \sum_j \alpha_j \ge 2$.
  Hence $f$ is a linear combination of $1, c_1, \ldots, c_m$, and $\varepsilon
  (f) = 0$ kills the constant.
\end{proof}

\subsection{The symmetric-function isomorphism}

Let $\Sym = \C [p_1, p_2, \ldots]$ be the Hopf algebra of symmetric functions
over $\C$, with power sums $p_r$ primitive, complete homogeneous symmetric
functions $h_N$, and Schur functions $s_{\lambda}$ {\tmem{(Macdonald
{\cite[Ch.~I]{Macdonald}})}}. Write $\CA_{\infty} \assign \C [c_1, c_2,
\ldots]$ for the evident $m = \infty$ version of Definitions~\ref{def:E}
and~\ref{def:hopf}; Lemmas \ref{lem:recurrence}--\ref{lem:derivation} and
Theorems~\ref{thm:generation}--\ref{thm:primitives} hold verbatim with $\min
(m, N)$ replaced by $N$.

\begin{theorem}
  [Symmetric-function realization]\label{thm:sym} The assignment $\Psi (p_j)
  \assign j \hspace{0.17em} c_j$ extends uniquely to an isomorphism of graded
  Hopf algebras $\Psi : \Sym \xrightarrow{\sim} \CA_{\infty}$, and
  \[ \Psi (h_N) = \Lambda_N  \qquad (N \ge 0) . \]
  Under $\Psi$, Newton's identity $Nh_N = \sum_{j = 1}^N p_j h_{N - j}$
  becomes the universal recurrence \eqref{eq:universal}. Moreover the ideal
  $I_m \assign (c_{m + 1}, c_{m + 2}, \ldots) \subset \CA_{\infty}$ is a
  graded Hopf ideal, and
  \[ \CA_m \hspace{0.27em} \cong \hspace{0.27em} \CA_{\infty} / I_m
     \hspace{0.27em} \cong \hspace{0.27em} \Sym / (p_{m + 1}, p_{m + 2},
     \ldots) \]
  as graded Hopf algebras, compatibly with the sequences: $h_N \mapsto
  \Lambda_N$ throughout.
\end{theorem}

\begin{proof}
  $\Sym$ and $\CA_{\infty}$ are free commutative algebras on primitive
  generators $p_j$, resp. $c_j$, one in each weight $j \ge 1$; $\Psi$ matches
  them by weight with the nonzero scalar $j$, hence is an isomorphism of
  graded algebras. It is a coalgebra map because it sends primitives to
  primitives and both coproducts are algebra morphisms determined on
  generators; a bialgebra isomorphism of Hopf algebras automatically respects
  antipodes. For the sequence: the generating identity $\sum_{N \ge 0} h_N t^N
  = \exp \left( \sum_{r \ge 1} p_r t^r / r \right)$
  {\tmem{({\cite[I.(2.10)]{Macdonald}})}} maps under $\Psi$ (applied
  coefficient-wise in $t$) to $\sum_N \Psi (h_N) t^N = \exp \left( \sum_r c_r
  t^r \right) = \sum_N \Lambda_N t^N$, so $\Psi (h_N) = \Lambda_N$. Newton's
  identity is {\cite[I.(2.11)$'$]{Macdonald}}; applying $\Psi$ gives
  \eqref{eq:universal} (and, independently, Lemma~\ref{lem:recurrence}
  re-proves Newton's identity via the differential equation $E' = (\sum jc_j
  v^{j - 1}) E$). For the quotient: $I_m$ is generated by primitives, so
  $\Delta I_m \subseteq I_m \otimes \CA_{\infty} + \CA_{\infty} \otimes I_m$,
  $\varepsilon (I_m) = 0$, $S (I_m) \subseteq I_m$: a Hopf ideal. The quotient
  map $\CA_{\infty} \to \CA_m$, $c_j \mapsto c_j$ $(j \le m)$, $c_j \mapsto 0$
  $(j > m)$, carries $\Lambda_N^{(\infty)}$ to $\Lambda_N$ by
  Lemma~\ref{lem:partition} (multi-indices with $k_j > 0$ for some $j > m$
  die).
\end{proof}

\begin{remark}
  \label{rem:notFdB}$\CA_{\infty}$ is {\tmem{not}} the Fa{\`a} di Bruno Hopf
  algebra: the latter is the coordinate Hopf algebra of the group of formal
  diffeomorphisms and is not cocommutative
  {\tmem{(Figueroa--Gracia-Bond{\'\i}a {\cite{FGB}})}}, whereas $\CA_{\infty}
  \cong \Sym$ is commutative and cocommutative. The relation to Bell
  polynomials is the coefficient identity of Proposition~\ref{prop:bell}, not
  an isomorphism with the Fa{\`a} di Bruno object. This answers, precisely,
  the question ``is it the Bell algebra'': the coefficient algebra is the
  algebra of symmetric functions, in which the $\Lambda_N$ are the complete
  homogeneous basis re-expressed through rescaled power sums.
\end{remark}

\begin{remark}
  [Hall algebra]\label{rem:hall} The classical Hall algebra of finite modules
  over a discrete valuation ring with finite residue field is isomorphic,
  after the standard rescaling, to $\Sym$
  {\tmem{({\cite[Ch.~II--III]{Macdonald}})}}. Via Theorem~\ref{thm:sym},
  $\CA_{\infty}$ therefore {\tmem{does}} arise as a Hall algebra, and $\CA_m$
  as an explicit Hopf quotient of one. We record this as an identification
  available by citation; we make no use of the finer (Hall--Littlewood, $q$-)
  structure here, and lifting the coefficient sequence $(\Lambda_N)$ to the
  $q$-deformed level is Open Problem~\ref{prob:hall}.
\end{remark}

\section{Combinatorial and matrix realizations}\label{sec:combinatorics}

\begin{theorem}
  [Species / set-partition interpretation]\label{thm:species} For $N \ge 1$,
  \[ N! \hspace{0.17em} \Lambda_N \hspace{0.27em} = \hspace{0.27em}
     \sum_{\tmscript{\begin{array}{c}
       \pi \vdash [N]\\
       \text{all blocks of size} \le m
     \end{array}}} \hspace{0.27em} \prod_{B \in \pi} |B| ! \hspace{0.27em}
     c_{|B|}, \]
  the sum over set partitions $\pi$ of $\{1, \ldots, N\}$ with block sizes at
  most $m$. Equivalently, $E (v)$ is the exponential generating function of
  the species ``sets of blocks,'' each block of size $j \le m$ weighted $j!
  \hspace{0.17em} c_j$.
\end{theorem}

\begin{proof}
  The number of set partitions of $[N]$ with exactly $k_j$ blocks of size $j$
  is $N! / \prod_j ((j!)^{k_j} k_j !)$. Hence
  \[ \sum_{\pi} \prod_{B \in \pi} |B| ! \hspace{0.17em} c_{|B|} =
     \sum_{\tmscript{\begin{array}{c}
       k : \hspace{0.17em} \sum jk_j = N
     \end{array}}} \frac{N!}{\prod_j (j!)^{k_j} k_j !}  \prod_j (j!
     \hspace{0.17em} c_j)^{k_j} = N! \sum_{\tmscript{\begin{array}{c}
       k : \hspace{0.17em} \sum jk_j = N
     \end{array}}} \prod_j \frac{c_j^{k_j}}{k_j !} = N! \hspace{0.17em}
     \Lambda_N \]
  by Lemma~\ref{lem:partition}. The species statement is the exponential
  formula {\tmem{(Flajolet--Sedgewick {\cite[II.2]{FS}};
  Bergeron--Labelle--Leroux {\cite{BLL}})}}.
\end{proof}

\begin{theorem}
  [Appell structure]\label{thm:appell} Fix $c_2, \ldots, c_m$ and regard $p_N
  (c_1) \assign N! \hspace{0.17em} \Lambda_N$ as a polynomial in $c_1$. Then
  $(p_N)_{N \ge 0}$ is an Appell sequence: $p_N' (c_1) = N \hspace{0.17em}
  p_{N - 1} (c_1)$, with exponential generating function $\sum_N p_N (c_1) t^N
  / N! = e^{c_1 t} A (t)$, $A (t) = \exp \left( \sum_{j = 2}^m c_j t^j
  \right)$.
\end{theorem}

\begin{proof}
  Lemma~\ref{lem:derivation} with $j = 1$ gives $\partial_{c_1} \Lambda_N =
  \Lambda_{N - 1}$, i.e. $p_N' = Np_{N - 1}$; the generating function is
  Definition~\ref{def:E} with the $c_1$-factor split off. Appell sequences are
  exactly the Sheffer sequences with identity delta series, so $(p_N)$ is
  Sheffer of Appell type {\tmem{(Roman {\cite[{\textsection}2.5]{Roman}})}}.
\end{proof}

\begin{theorem}
  [Riordan-group realization]\label{thm:riordan} The lower-triangular matrix
  $D = (d_{n, k})_{n, k \ge 0}$, $d_{n, k} \assign \dfrac{n!}{k!}
  \hspace{0.17em} \Lambda_{n - k}$, is the exponential Riordan array $\left( E
  (t), \hspace{0.17em} t \right)$, an element of the Appell subgroup of the
  exponential Riordan group. The group law restricted to this subgroup, $(E
  (t), t) \cdot (E' (t), t) = (E (t) E' (t), t)$, corresponds to addition $c_j
  \mapsto c_j + c_j'$ of parameters, i.e. to the coproduct of
  Theorem~\ref{thm:dividedpower}.
\end{theorem}

\begin{proof}
  By definition {\tmem{(Shapiro et al. {\cite{Shapiro}}; Barry
  {\cite{Barry}})}} the exponential Riordan array $(g (t), f (t))$ has entries
  $d_{n, k} = \frac{n!}{k!} [t^{n - k}]  \hspace{0.17em} g (t) f (t)^k / t^k$
  reduced, for $f (t) = t$, to $d_{n, k} = \frac{n!}{k!} [t^{n - k}] g (t)$;
  with $g = E$ this is $\frac{n!}{k!} \Lambda_{n - k}$. Matrix multiplication
  of Appell arrays multiplies the $g$'s; $E (t ; c) E (t ; c') = E (t ; c +
  c')$ since the exponents add. The identification with the coproduct is
  Theorem~\ref{thm:dividedpower} read through $\Lambda_N  (c + c') = \sum_i
  \Lambda_i (c) \Lambda_{N - i} (c')$, which is the image of $\Delta
  \Lambda_N$ under evaluation at $(c, c')$.
\end{proof}

\subsection{Hankel determinants are rectangular Schur functions}

\begin{theorem}
  [Hankel--Schur identity]\label{thm:hankelschur} In $\CA_{\infty} \cong \Sym$
  (hence, after applying the quotient of Theorem~\ref{thm:sym}, as identities
  in $\CA_m$), for every $n \ge 1$:
  \[ \det (\Lambda_{i + j})_{0 \le i, j \le n - 1} = (- 1)^{\binom{n}{2}} 
     \hspace{0.17em} \Psi (s_{((n - 1)^n)}), \qquad \det (\Lambda_{i + j +
     1})_{0 \le i, j \le n - 1} = (- 1)^{\binom{n}{2}}  \hspace{0.17em} \Psi
     (s_{(n^n)}), \]
  where $(r^n)$ denotes the $n \times r$ rectangular partition.
\end{theorem}

\begin{proof}
  Work in $\Sym$ with $\Lambda_N = h_N$ (Theorem~\ref{thm:sym}). Reverse the
  rows of the Hankel matrix by $i \mapsto n - 1 - i$; this multiplies the
  determinant by the sign of the reversal permutation, $(- 1)^{\binom{n}{2}}$.
  In $1$-based indices $(i, j)$ the reversed matrix has entries
  $h_{\hspace{0.17em} n - 1 - i + j}$ (first case), which is the Jacobi--Trudi
  matrix $(h_{\lambda_i - i + j})$ for the constant partition $\lambda_i = n -
  1$; by the Jacobi--Trudi identity {\tmem{({\cite[I.(3.4)]{Macdonald}})}} its
  determinant is $s_{((n - 1)^n)}$. The shifted case is identical with
  $\lambda_i = n$.
\end{proof}

\begin{corollary}
  [Normality of the Pad{\'e} table of $E$]\label{cor:padeE} Let $\ev_c : \CA_m
  \to \C$ be evaluation at $c \in \C^m$ and let $\mathcal{L}$ be the moment
  functional $\mathcal{L} [x^N] = \ev_c (\Lambda_N)$. The Hankel determinants
  governing the diagonal Pad{\'e} approximants of $E (v)$ at $c$ {\tmem{(Gragg
  {\cite{Gragg}})}} are the numbers $(- 1)^{\binom{n}{2}}  \hspace{0.17em}
  s_{((n - 1)^n)}$ and $(- 1)^{\binom{n}{2}}  \hspace{0.17em} s_{(n^n)}$
  evaluated under $p_j = jc_j$ $(j \le m)$, $p_j = 0$ $(j > m)$. In
  particular, for $m = 1$ (so $p_1 = c_1$, $p_{\ge 2} = 0$), the exponential
  specialization $s_{\lambda} \mapsto f^{\lambda} c_1^{| \lambda |} / |
  \lambda |$! {\tmem{({\cite[I.{\textsection}7, Ex.~24(a) with $n \to \infty$;
  equivalently $\ev$ of {\cite[I.(5.16)]{Macdonald}} at the exponential
  specialization]{Macdonald}}{\hspace{0.17em}})}}, where $f^{\lambda} > 0$ is
  the number of standard Young tableaux, shows every such determinant is
  nonzero for $c_1 \ne 0$: the Pad{\'e} table of $e^{c_1 v}$ is normal,
  recovering the classical normality of the exponential.
\end{corollary}

\begin{proof}
  Immediate from Theorem~\ref{thm:hankelschur}, the Pad{\'e}--Hankel
  dictionary {\cite{Gragg}}, and $f^{\lambda} \ge 1$ for every partition
  $\lambda$.
\end{proof}

\begin{remark}
  \label{rem:zeroloci}For $m \ge 2$ the quasi-definiteness locus of
  $\mathcal{L}$ is precisely the complement of the union of the hypersurfaces
  $\{c \in \C^m : s_{((n - 1)^n)} (p_j = jc_j) = 0\}$; the obstruction to
  Pad{\'e} normality of exponentials of principal parts is a rectangular Schur
  divisor. This gives the ``discriminant locus'' appearing in the companion
  Chebyshev--Wheeler paper a name in the pure-pole case.
\end{remark}

\section{Automorphisms and differential Galois structure}\label{sec:aut}

\begin{theorem}
  [Graded Hopf automorphisms]\label{thm:aut} $\Aut_{\textrm{\tmop{gr} \text{-}
  \tmop{Hopf}}} (\CA_m) \cong (\C^{\times})^m$, acting by $c_j \mapsto
  \lambda_j c_j$. The subgroup induced by the geometric rescalings $w \mapsto
  \lambda w$ of the local coordinate is the one-parameter diagonal $\lambda_j
  = \lambda^{\hspace{0.17em} j}$.
\end{theorem}

\begin{proof}
  A graded Hopf automorphism preserves $\Prim (\CA_m)$ and each of its graded
  pieces, which are the lines $\C c_j$ by Theorem~\ref{thm:primitives}; hence
  $c_j \mapsto \lambda_j c_j$ with $\lambda_j \in \C^{\times}$, and any such
  assignment plainly defines a graded Hopf automorphism. Under $w \mapsto
  \lambda w$ the principal part $\sum c_j w^{- j}$ becomes $\sum (c_j
  \lambda^{- j}) w^{- j}$, giving the stated one-parameter subgroup (with
  $\lambda^{- j}$; reparametrize).
\end{proof}

\begin{proposition}
  [Picard--Vessiot group]\label{prop:galois} Let $R \in \C (z)$ be
  nonconstant. Then $y = e^R$ is transcendental over $\C (z)$ and the
  Picard--Vessiot group of $y' = R' y$ over $\C (z)$ is $\mathbb{G}_m$. The
  group acts on solutions by $y \mapsto \lambda y$, hence on the full Laurent
  coefficient sequence $(A_n)$ of Section~\ref{sec:model} diagonally by $A_n
  \mapsto \lambda A_n$: the coefficient sequence spans the weight-one isotypic
  line, while the algebra $\CA_m$ itself is Galois-invariant.
\end{proposition}

\begin{proof}
  The PV group of a rank-one equation is a Zariski-closed subgroup of
  $\mathbb{G}_m$, hence $\mathbb{G}_m$ itself or the finite group $\mu_n$; the
  latter occurs iff $y^n \in \C (z)$ for some $n \ge 1$ {\tmem{(van der
  Put--Singer {\cite[Ex.~1.19, {\textsection}1.4]{vdPS}})}}. Suppose $e^{nR} =
  f \in \C (z)$. Then $nR' = f' / f$. Every residue of $R'$ vanishes, since
  $R'$ is the derivative of a rational function; every residue of $f' / f$ at
  a zero or pole of $f$ is the (nonzero) integer order of $f$ there. Hence $f$
  has no zeros or poles, so $f$ is constant and $R' = 0$, contradicting $R$
  nonconstant.
\end{proof}

\section{The analytic model}\label{sec:model}

We now verify that the formal object of
Sections~\ref{sec:recurrence}--\ref{sec:combinatorics} computes actual Laurent
coefficients, with quantitative control. Fix $c = (c_1, \ldots, c_m) \in \C^m$
and write $\Lambda_N (c) \assign \ev_c (\Lambda_N)$, $C \assign \sum_{j = 1}^m
|c_j |$.

\begin{proposition}
  [Growth bound]\label{prop:growth} For all $N \ge 0$ and all $v > 0$, $|
  \Lambda_N (c) | \le \Lambda_N  (|c_1 |, \ldots, |c_m |) \le v^{- N} \exp
  \left( \sum_j |c_j |v^j \right)$. Consequently, for $N \ge mC$,
  \[ | \Lambda_N (c) | \hspace{0.27em} \le \hspace{0.27em} e^{N / m} \left(
     \frac{mC}{N} \right)^{N / m}, \]
  so $| \Lambda_N (c) |^{1 / N} \to 0$ super-geometrically, at the rate of an
  entire function of order $m$ in $1 / w$.
\end{proposition}

\begin{proof}
  The first inequality is termwise from \eqref{eq:partition}. For the second,
  all terms of $\sum_M \Lambda_M (|c|) v^M = \exp (\sum_j |c_j |v^j)$ are
  nonnegative, so any single term is bounded by the sum: $\Lambda_N (|c|) v^N
  \le \exp (\sum |c_j |v^j) \le e^{Cv^m}$ for $v \ge 1$. Choose $v = (N /
  (mC))^{1 / m} \ge 1$; then $e^{Cv^m} = e^{N / m}$ and $v^{- N} = (mC / N)^{N
  / m}$.
\end{proof}

\begin{theorem}
  [Pure-pole expansion]\label{thm:analyticmodel} For every $c \in \C^m$ and $w
  \ne 0$, the series $\sum_{N \ge 0} \Lambda_N (c)  \hspace{0.17em} w^{- N}$
  converges absolutely, uniformly on $|w| \ge \epsilon$ for each $\epsilon >
  0$, and
  \[ \exp \left( \sum_{j = 1}^m c_j w^{- j} \right) = \sum_{N \ge 0} \Lambda_N
     (c)  \hspace{0.17em} w^{- N} . \]
  In particular $\Lambda_N (c)$ {\tmem{is}} the coefficient of $w^{- N}$ in
  the Laurent expansion of $e^{\Pi}$ on $\C^{\times}$, and the map $c \mapsto
  (\Lambda_N (c))_{N \ge 0}$ is injective.
\end{theorem}

\begin{proof}
  Expanding each factor $\exp (c_j w^{- j})$ and multiplying, the resulting
  multi-series is dominated termwise by $\prod_j \sum_{k_j} |c_j |^{k_j}
  |w|^{- jk_j} / k_j ! = \exp (\sum_j |c_j ||w|^{- j}) < \infty$, so it
  converges absolutely and may be rearranged by total degree, yielding $\sum_N
  \Lambda_N (c) w^{- N}$ with the coefficients of Lemma~\ref{lem:partition}.
  Uniqueness of Laurent coefficients on an annulus identifies them with the
  Laurent coefficients of $e^{\Pi}$. Injectivity: by the triangularity in the
  proof of Theorem~\ref{thm:generation}, $c_N$ is a polynomial in $\Lambda_1
  (c), \ldots, \Lambda_N (c)$.
\end{proof}

\begin{theorem}
  [Local Laurent structure of $e^R$; the coefficient module]\label{thm:module}
  Let $R \in \C (z)$ have a pole at $a$ of order $m \ge 1$ with principal part
  $\Pi (w) = \sum_{j = 1}^m c_j w^{- j}$, $w = z - a$, $c_m \ne 0$, and let
  $r_a \in (0, \infty]$ be the distance from $a$ to the nearest other pole of
  $R$. Write $H \assign R - \Pi$, analytic on $|w| < r_a$, and $e^{H (a + w)}
  = \sum_{l \ge 0} b_l w^l$, $\limsup_l |b_l |^{1 / l} \le 1 / r_a$. Then on
  $0 < |w| < r_a$,
  \[ e^{R (a + w)} = \sum_{n \in \Z} A_n  \hspace{0.17em} w^n, \qquad A_n =
     \sum_{k \ge \max (0, - n)} \Lambda_k (c)  \hspace{0.17em} b_{n + k}, \]
  each $A_n$ given by an absolutely convergent series. The map $\CA_m \times
  \C [[w]]_{r_a} \to \C^{\Z}$ so defined is $\C$-bilinear in the coefficient
  sequences $(\Lambda_k (c))$ and $(b_l)$: the Laurent coefficients of $e^R$
  constitute a module over the coefficient algebra, with structure constants
  supported on the exponent lattice $n = l - k$.
\end{theorem}

\begin{proof}
  $e^R = e^{\Pi} e^H$ exactly, since $\Pi$ and $H$ are scalar functions. On
  any circle $|w| = \rho \in (0, r_a)$, both factors are given by absolutely
  convergent Laurent/Taylor series (Theorem~\ref{thm:analyticmodel}; Cauchy),
  so their product's Laurent coefficients are the Cauchy convolutions $A_n =
  \sum_k \Lambda_k (c) b_{n + k}$ (with $b_{< 0} \assign 0$), and the double
  series converges absolutely because, choosing $\rho' \in (\rho, r_a)$,
  $|b_{n + k} | {\le B \rho'}^{- (n + k)}$ while $| \Lambda_k (c) |$ decays
  super-geometrically by Proposition~\ref{prop:growth}: $\sum_k | \Lambda_k
  (c) |  \hspace{0.17em} {B \rho'}^{- (n + k)} {\le B \rho'}^{- n}  \sum_k
  e^{k / m}  (mC / k)^{k / m} {\rho'}^{- k} < \infty$. Bilinearity is
  manifest.
\end{proof}

\begin{proposition}
  [Truncation error]\label{prop:tail} In the setting of
  Theorem~\ref{thm:module}, fix $\rho' < r_a$ with $|b_l | {\le B \rho'}^{-
  l}$, and let $K \ge K_0 \assign \max \left\{ mC, e \hspace{0.17em} mC
  \hspace{0.17em} (2 / \rho')^m \right\}$. Then
  \[ \left| A_n - \sum_{k = \max (0, - n)}^K \Lambda_k (c) b_{n + k} \right|
     \hspace{0.27em} \le \hspace{0.27em} B \hspace{0.17em} {\rho'}^{- n} 
     \hspace{0.17em} 2^{- K} . \]
\end{proposition}

\begin{proof}
  For $k \ge K_0$, Proposition~\ref{prop:growth} gives $| \Lambda_k (c) |
  {\rho'}^{- k} \le \left[ (e \hspace{0.17em} mC / k)^{1 / m} / \rho'
  \right]^k \le 2^{- k}$ by the choice of $K_0$. Sum the geometric tail
  $\sum_{k > K} 2^{- k} \le 2^{- K}$ against ${B \rho'}^{- n}$ (absorbing
  ${\rho'}^{- k} \le$ what was already used).
\end{proof}

\begin{proposition}
  [Hadamard closure]\label{prop:hadamard} For $c, c' \in \C^m$, the Hadamard
  product sequence $(\Lambda_N (c) \Lambda_N (c'))_N$ is P-recursive
  (holonomic).
\end{proposition}

\begin{proof}
  Each factor satisfies the polynomial-coefficient recurrence
  \eqref{eq:universal}; the Hadamard product of P-recursive sequences is
  P-recursive {\tmem{(Stanley {\cite{StanleyDF}}; Lipshitz
  {\cite{Lipshitz}})}}.
\end{proof}

\section{Classical specializations}\label{sec:classical}

The identifications below are exact and exhaustive in the stated cases; no
claim is made that a general principal part of order $m \ge 3$ reduces to a
single classical one-variable function. The general case is the algebra
itself.

\begin{theorem}
  [$m = 1$: modified Bessel]\label{thm:bessel} For $b, c \in \C$, the Laurent
  expansion of $e^{c / w + bw}$ on $\C^{\times}$ has coefficients
  \[ [w^{- n}]  \hspace{0.17em} e^{c / w + bw} = \sum_{j \ge 0}
     \frac{c^{\hspace{0.17em} j + n} b^{\hspace{0.17em} j}}{j! \hspace{0.17em}
     (j + n) !}  \qquad (n \ge 0), \qquad [w^n]  \hspace{0.17em} e^{c / w +
     bw} = \text{same with } b \leftrightarrow c. \]
  For any determination $\sqrt{bc}$ of a square root of $bc$ with
  $(\sqrt{bc})^2 = bc$, these equal $(c / b)^{n / 2} I_n  (2 \sqrt{bc})$
  (interpreting $(c / b)^{n / 2} \assign c^n  (bc)^{- n / 2}$ with the same
  determination), where $I_n (x) = \sum_{j \ge 0} (x / 2)^{2 j + n} / (j!
  \hspace{0.17em} (j + n) !)$. Setting $b = c = x / 2$, $w = 1 / t$ recovers
  the classical generating identity $\exp \left( \tfrac{x}{2} (t + 1 / t)
  \right) = \sum_{n \in \Z} I_n (x) t^n$.
\end{theorem}

\begin{proof}
  Here $\Lambda_k (c) = c^k / k$! and $b_l = b^l / l$!;
  Theorem~\ref{thm:module} with $r_a = \infty$ gives $A_{- n} = \sum_{k \ge n}
  \frac{c^k}{k!}  \frac{b^{k - n}}{(k - n) !}$; substitute $k = j + n$. The
  Bessel form is the same series regrouped: $\sum_j (bc)^j / (j! (j + n) !) =
  (bc)^{- n / 2} I_n  (2 \sqrt{bc})$ term by term, both sides being
  unambiguous power series in $bc$ once the prefactor convention is fixed.
  Symmetry $w \mapsto 1 / w$ exchanges $b, c$ and $n \mapsto - n$, giving the
  positive coefficients and, with $I_{- n} = I_n$, the classical identity
  {\tmem{(cf. {\cite[10.35.1]{DLMF}})}}.
\end{proof}

\begin{theorem}
  [$m = 2$: Hermite-type polynomials]\label{thm:hermite} Define $\mathbf{H}_N
  (x, y)$ by $\exp (xt + yt^2) = \sum_{N \ge 0} \mathbf{H}_N (x, y) 
  \hspace{0.17em} t^N / N!$ {\tmem{(the Hermite--Kamp{\'e} de F{\'e}riet
  polynomials {\cite{Appell}})}}. Then for $m = 2$,
  \[ \Lambda_N  (c_1, c_2) = \frac{\mathbf{H}_N (c_1, c_2)}{N!} = \sum_{0 \le
     k \le N / 2} \frac{c_1^{\hspace{0.17em} N - 2 k} c_2^{\hspace{0.17em}
     k}}{(N - 2 k) ! \hspace{0.17em} k!}, \]
  and $\mathbf{H}_N (x, y) = (\mathrm{i} \sqrt{y})^{\hspace{0.17em} N} H_N 
  \hspace{-0.17em} \left( x / (2 \mathrm{i} \sqrt{y}) \right)$ for any fixed
  square root, where $H_N$ is the classical (physicists') Hermite polynomial.
\end{theorem}

\begin{proof}
  The first two equalities are Definition~\ref{def:E} for $m = 2$ and
  Lemma~\ref{lem:partition}. The classical relation follows by comparing $\exp
  (xt + yt^2)$ with $\exp (2 Xs - s^2) = \sum H_N (X) s^N / N!$ under $s =
  \mathrm{i} \sqrt{y}  \hspace{0.17em} t$, $X = x / (2 \mathrm{i} \sqrt{y})$;
  both sides are polynomial identities in $(x, y^{1 / 2})$, hence hold for
  every branch choice consistently.
\end{proof}

\begin{theorem}
  [Single-term pole of order $m$: Wright functions]\label{thm:wright} Let
  $\phi (\rho, \beta ; \zeta) \assign \sum_{k \ge 0} \dfrac{\zeta^k}{k!
  \hspace{0.17em} \Gamma (\rho k + \beta)}$ be Wright's generalized Bessel
  function {\tmem{(Wright {\cite{Wright1933}}; Kilbas--Saigo
  {\cite[{\textsection}1.11]{KilbasSaigo}})}}, entire in $\zeta$ for $\rho > -
  1$. Then for $b \in \C^{\times}$, $c \in \C$, $n \in \Z$,
  \[ [w^{- n}]  \hspace{0.27em} \exp \left( \frac{c}{w^m} + bw \right) = b^{-
     n}  \hspace{0.17em} \phi \left( m, \hspace{0.17em} 1 - n ;
     \hspace{0.17em} c \hspace{0.17em} b^{\hspace{0.17em} m} \right) . \]
\end{theorem}

\begin{proof}
  Here $\Lambda_{mk} = c^k / k$! and $\Lambda_N = 0$ otherwise, $b_l = b^l /
  l$! $(l \ge 0)$, $b_l \assign 0$ $(l < 0)$. Theorem~\ref{thm:module}: $A_{-
  n} = \sum_k \frac{c^k}{k!}  \hspace{0.17em} \frac{b^{\hspace{0.17em} mk -
  n}}{(mk - n) !}$ over $k$ with $mk \ge n$. Since $1 / \Gamma (\ell) = 0$ for
  $\ell \in \Z_{\le 0}$, the constraint is absorbed by writing $(mk - n) ! =
  \Gamma (mk - n + 1) = \Gamma (mk + (1 - n))$, and the full sum over $k \ge
  0$ equals $b^{- n}  \sum_k (cb^m)^k / \left( k! \hspace{0.17em} \Gamma (mk +
  1 - n) \right)$.
\end{proof}

\begin{remark}
  [Scope of special-function reductions]\label{rem:scope}
  Theorems~\ref{thm:bessel}--\ref{thm:wright} exhaust the cases in which a
  single classical family arises from a natural one-parameter entire factor:
  $m = 1$, $m = 2$, and pure top-order poles. For a general principal part
  with $m \ge 3$ and mixed $c_j$, the coefficient $A_n$ is the genuine
  multivariate partition sum of Theorem~\ref{thm:module}, i.e. an element of
  the $\CA_m$-module, and we do not assert any reduction to Wright,
  Fox--Wright, or Fox $H$ functions of one variable. What replaces such a
  reduction is the structure theory of
  Sections~\ref{sec:hopf}--\ref{sec:combinatorics} together with the holonomy
  of Section~\ref{sec:holonomy}: closed recurrences, not closed one-variable
  forms.
\end{remark}

\section{Holonomy}\label{sec:holonomy}

\begin{theorem}
  [D-finiteness and the coefficient recurrence]\label{thm:holonomy} Let $R \in
  \C (z)$ be nonconstant, write $R' = U / V$ in lowest terms with $U, V \in \C
  [z]$, and let $a$ be a pole of $R$ of order $m \ge 1$. Set $\tilde{V} (w)
  \assign V (a + w) = \sum_{p = 0}^{d_V} V_p w^p$, $\tilde{U} (w) \assign U (a
  + w) = \sum_{p = 0}^{d_U} U_p w^p$. Then $y = e^R$ satisfies the first-order
  equation $Vy' - Uy = 0$ over $\C (z)$ (in particular $y$ is D-finite), and
  the full two-sided Laurent coefficient sequence $(A_n)_{n \in \Z}$ of $y$ at
  $a$ (Theorem~\ref{thm:module}) satisfies, for every $s \in \Z$,
  \begin{equation}
    \label{eq:Prec} \sum_{p = 0}^{d_V} V_p  \hspace{0.17em} (s + 1 - p) 
    \hspace{0.17em} A_{s + 1 - p} \hspace{0.27em} = \hspace{0.27em} \sum_{p =
    0}^{d_U} U_p  \hspace{0.17em} A_{s - p} .
  \end{equation}
  This is a linear recurrence with coefficients affine in $s$, of order at
  most $\max (d_V, \hspace{0.17em} d_U + 1)$, valid in both directions $n \to
  \pm \infty$. Moreover $w^{m + 1} \mid \tilde{V}$ (indeed $\ord_{w = 0} R' =
  - (m + 1)$ forces $\ord \tilde{V} - \ord \tilde{U} = m + 1$), and in the
  pure-pole case $R = \Pi$ the recurrence \eqref{eq:Prec} reduces exactly to
  the universal recurrence \eqref{eq:universal}.
\end{theorem}

\begin{proof}
  $y' = R' y$ is immediate; clearing denominators gives $Vy' = Uy$ with
  polynomial coefficients, which is the definition of D-finiteness for $y$
  {\tmem{(Stanley {\cite{StanleyDF}})}}. On $0 < |w| < r_a$ insert $y = \sum_n
  A_n w^n$ (absolutely and locally uniformly convergent, so termwise
  differentiation and Cauchy products with the polynomials $\tilde{V},
  \tilde{U}$ are valid): $\sum_p V_p w^p  \sum_n nA_n w^{n - 1} = \sum_p U_p
  w^p  \sum_n A_n w^n$. Equate coefficients of $w^s$: $\sum_p V_p  (s + 1 - p)
  A_{s + 1 - p} = \sum_p U_p A_{s - p}$, which is \eqref{eq:Prec}. The order
  statement counts the index span. The valuation claim: near $a$, $R' = - mc_m
  w^{- (m + 1)} + O (w^{- m})$ with $c_m \ne 0$, so $\ord_{w = 0} (\tilde{U} /
  \tilde{V}) = - (m + 1)$; since $\gcd (U, V) = 1$, at most one of $\tilde{U}
  (0), \tilde{V} (0)$ vanishes, forcing $\tilde{U} (0) \ne 0$ and $w^{m + 1}
  \hspace{0.17em} \| \hspace{0.17em} \tilde{V}$. Pure pole: $R = \Pi$ gives
  $R' = - \sum_j jc_j w^{- j - 1} = \tilde{U} / \tilde{V}$ with $\tilde{V} =
  w^{m + 1}$, $\tilde{U} = - \sum_{j = 1}^m jc_j w^{\hspace{0.17em} m - j}$;
  then \eqref{eq:Prec} reads $(s - m) A_{s - m} = - \sum_j jc_j A_{s - m +
  j}$, and putting $A_{- N} = \Lambda_N$, $s - m = - N$: $N \Lambda_N = \sum_j
  jc_j \Lambda_{N - j}$.
\end{proof}

\begin{corollary}
  \label{cor:algorithmicholonomy}For fixed $c$, every coefficient sequence in
  the module of Theorem~\ref{thm:module} is P-recursive with an explicitly
  computable recurrence \eqref{eq:Prec}; creative-telescoping toolchains apply
  directly {\tmem{(Kauers--Paule {\cite{KauersPaule}}; Koutschan
  {\cite{Koutschan}})}}.
\end{corollary}

\section{Axiomatic characterization}\label{sec:axioms}

Following the constructive development, we now characterize what has been
built. Two statements are proved: an {\tmem{initiality}} theorem in the
category naturally attached to the universal recurrence, and a
{\tmem{recognition}} theorem identifying $\CA_m$ among graded Hopf algebras.
The hypotheses of the recognition theorem are shown to be independent.

\subsection{Initiality}

\begin{definition}
  [The category $\mathcal{C}_m$]\label{def:category} Objects of
  $\mathcal{C}_m$ are triples $(A, \gamma = (\gamma_1, \ldots, \gamma_m) \in
  A^m, (\lambda_N)_{N \ge 0})$ where $A$ is a commutative unital $\C$-algebra,
  $\lambda_0 = 1$, and
  \begin{equation}
    \label{eq:axiomrec} N \lambda_N = \sum_{j = 1}^{\min (m, N)} j
    \hspace{0.17em} \gamma_j  \hspace{0.17em} \lambda_{N - j}  \qquad (N \ge
    1) .
  \end{equation}
  A morphism $(A, \gamma, \lambda) \to (A', \gamma', \lambda')$ is a unital
  $\C$-algebra homomorphism $\varphi : A \to A'$ with $\varphi (\gamma_j) =
  \gamma'_j$ for all $j$ and $\varphi (\lambda_N) = \lambda'_N$ for all $N$.
\end{definition}

\begin{theorem}
  [Initiality]\label{thm:initial} $(\CA_m, \hspace{0.17em} c, \hspace{0.17em}
  \Lambda)$ is an initial object of $\mathcal{C}_m$: for every object $(A,
  \gamma, \lambda)$ there is exactly one morphism $\varphi : (\CA_m, c,
  \Lambda) \to (A, \gamma, \lambda)$.
\end{theorem}

\begin{proof}
  {\tmem{Existence.}} Since $\CA_m$ is the free commutative $\C$-algebra on
  $c_1, \ldots, c_m$, there is a unique algebra map $\varphi$ with $\varphi
  (c_j) = \gamma_j$. We claim $\varphi (\Lambda_N) = \lambda_N$ for all $N$,
  by strong induction: true at $N = 0$ ($\Lambda_0 = 1 \mapsto 1 =
  \lambda_0$); for $N \ge 1$, since $N$ is invertible in the $\C$-algebra $A$,
  \[ \varphi (\Lambda_N) = \varphi \left( \tfrac{1}{N} \sum_j jc_j \Lambda_{N
     - j} \right) = \tfrac{1}{N} \sum_j j \hspace{0.17em} \gamma_j 
     \hspace{0.17em} \varphi (\Lambda_{N - j}) = \tfrac{1}{N} \sum_j j
     \hspace{0.17em} \gamma_j  \hspace{0.17em} \lambda_{N - j} = \lambda_N, \]
  using Lemma~\ref{lem:recurrence} in $\CA_m$, the inductive hypothesis, and
  \eqref{eq:axiomrec} in $A$. {\tmem{Uniqueness.}} Any morphism must send $c_j
  \mapsto \gamma_j$, and an algebra map out of a polynomial ring is determined
  by the images of the variables; the $\lambda$-compatibility is then
  automatic by the computation just made, so no further freedom exists.
\end{proof}

\begin{remark}
  \label{rem:initialmeaning}Equation \eqref{eq:axiomrec} determines
  $(\lambda_N)$ uniquely from $\gamma$ in any $\C$-algebra; so $\mathcal{C}_m$
  is equivalent to the category of $m$-pointed commutative $\C$-algebras, and
  initiality is exactly the freeness of $\CA_m$ {\tmem{together with}} the
  verification that the distinguished sequence transports along the unique
  map. The content is modest but the statement is now a theorem with a defined
  category, defined morphisms, and a proof, rather than an assertion.
\end{remark}

\subsection{Recognition}

\begin{theorem}
  [Recognition]\label{thm:recognition} Let $A$ be a graded connected
  commutative $\C$-Hopf algebra with a sequence $\lambda_N \in A_N$,
  $\lambda_0 = 1$, satisfying:
  \begin{enumerate}
    \item[\tmtextup{(G)}] the $\lambda_N$ generate $A$ as an algebra;
    
    \item[\tmtextup{(D)}] $\Delta \lambda_N = \sum_{i = 0}^N \lambda_i \otimes
    \lambda_{N - i}$ for all $N$ {\tmem{(divided-power sequence)}};
    
    \item[\tmtextup{(T)}] $\Prim (A)_j = 0$ for all $j > m$;
    
    \item[\tmtextup{(I)}] $\lambda_1, \ldots, \lambda_m$ are algebraically
    independent over $\C$.
  \end{enumerate}
  Then there is a unique isomorphism of graded Hopf algebras $\varphi : \CA_m
  \to A$ with $\varphi (\Lambda_N) = \lambda_N$ for all $N$; and $\varphi
  (c_j)$ is the weight-$j$ coefficient of $\log \sum_N \lambda_N t^N$.
\end{theorem}

\begin{proof}
  Set $\Lambda^A (t) \assign \sum_{N \ge 0} \lambda_N t^N \in A [[t]]$. By
  (D), $\Lambda^A (t)$ is group-like in the completed Hopf algebra $A [[t]]$
  (coproduct and counit extended $t$-adically): $\Delta \Lambda^A = \Lambda^A
  \otimes \Lambda^A$ read in $(A \otimes A) [[t]]$, $\varepsilon \Lambda^A =
  1$. Since $\Lambda^A (t) - 1 \in tA [[t]]$, the logarithm $\ell (t) \assign
  \log \Lambda^A (t) = \sum_{j \ge 1} \gamma_j t^j$ is defined, and
  \[ \Delta \ell = \log ((\Lambda^A \otimes 1) (1 \otimes \Lambda^A)) = \log
     (\Lambda^A \otimes 1) + \log (1 \otimes \Lambda^A) = \ell \otimes 1 + 1
     \otimes \ell, \]
  the middle equality because the two arguments commute. Hence every
  $\gamma_j$ is primitive; also $\gamma_j \in A_j$ because $\lambda_N \in A_N$
  and $\log$ preserves the grading of coefficients. By (T), $\gamma_j = 0$ for
  $j > m$, so
  \[ \Lambda^A (t) = \exp \left( \sum_{j = 1}^m \gamma_j t^j \right), \qquad
     \text{i.e.} \qquad \lambda_N = \Lambda_N  (\gamma_1, \ldots, \gamma_m) .
  \]
  Define the algebra map $\varphi : \CA_m \to A$ by $\varphi (c_j) =
  \gamma_j$; then $\varphi (\Lambda_N) = \lambda_N$ by the displayed identity
  (or by Theorem~\ref{thm:initial}, since $(A, \gamma, \lambda) \in
  \mathcal{C}_m$: apply $d / dt$ to $\Lambda^A = \exp (\sum \gamma_j t^j)$ as
  in Lemma~\ref{lem:recurrence}). $\varphi$ is a morphism of graded Hopf
  algebras: it preserves grading ($\gamma_j \in A_j$) and sends the primitive
  generators $c_j$ to primitives, which determines compatibility with
  $\Delta$, $\varepsilon$, $S$ on all of $\CA_m$ since these are (anti)algebra
  maps. {\tmem{Surjectivity:}} the image contains every $\lambda_N$, which
  generate by (G). {\tmem{Injectivity:}} $\lambda_1, \ldots, \lambda_m$ are
  polynomials in $\gamma_1, \ldots, \gamma_m$ and conversely (the
  triangularity of Theorem~\ref{thm:generation} transports along $\varphi$),
  so (I) is equivalent to algebraic independence of $\gamma_1, \ldots,
  \gamma_m$; an algebra map from $\C [c_1, \ldots, c_m]$ sending the variables
  to algebraically independent elements has zero kernel. {\tmem{Uniqueness:}}
  any graded Hopf map $\psi$ with $\psi (\Lambda_N) = \lambda_N$ satisfies
  $\psi (\log E (t)) = \ell (t)$ coefficient-wise, i.e. $\psi (c_j) =
  \gamma_j$, and is therefore $\varphi$.
\end{proof}

\subsection{Independence of the hypotheses}

\begin{proposition}
  \label{prop:independence}Each of \tmtextup{(G), (D), (T), (I)} is
  independent of the remaining three: for each hypothesis there is a triple
  $(A, (\lambda_N))$ satisfying the other three but not isomorphic to $(\CA_m,
  (\Lambda_N))$ in the sense of Theorem~\ref{thm:recognition}.
\end{proposition}

\begin{proof}
  Fix $m \ge 2$ (for $m = 1$ replace the (I)-example by $A = \C$, $\lambda_N =
  \delta_{N, 0}$, which satisfies (G),(D),(T) but not (I)).
  
  {\tmem{Drop (G).}} $A = \CA_m [y]$, $y$ primitive of weight $1$, with
  $\lambda_N \assign \Lambda_N$. (D) holds as before; $\Prim (A)_j = \C c_j$
  for $2 \le j \le m$, $\Prim (A)_1 = \C c_1 \oplus \C y$, and $0$ for $j >
  m$, so (T) holds; (I) holds. But the $\lambda_N$ do not generate ($y$ is
  missing), and $A \ncong \CA_m$ carrying the sequences (Hilbert series
  differ).
  
  {\tmem{Drop (I).}} $A = \CA_m / (c_m) \cong \CA_{m - 1}$ with $\lambda_N
  \assign$ image of $\Lambda_N$, i.e. $\Lambda_N  (c_1, \ldots, c_{m - 1},
  0)$. $(c_m)$ is a Hopf ideal (generated by a primitive), so $A$ is a graded
  connected Hopf algebra and (D) descends; (G) descends; (T) holds a fortiori
  since $\Prim (A)_j = 0$ for $j \ge m$. But $\lambda_1, \ldots, \lambda_m$
  satisfy the algebraic relation expressing $\lambda_m$ through $\lambda_1,
  \ldots, \lambda_{m - 1}$ (set $c_m = 0$ in $\Lambda_m = c_m + q_m (c_1,
  \ldots, c_{m - 1})$), and $A \ncong \CA_m$ (transcendence degree $m - 1$ vs.
  $m$).
  
  {\tmem{Drop (T).}} $A = \CA_{m + 1}$ with its own sequence $\lambda_N
  \assign \Lambda^{(m + 1)}_N$. (G), (D), (I) hold (Theorems
  \ref{thm:generation}, \ref{thm:dividedpower}; independence of the first $m$
  among the first $m + 1$). $\Prim (A)_{m + 1} = \C c_{m + 1} \ne 0$, and no
  isomorphism onto $\CA_m$ exists (transcendence degree).
  
  {\tmem{Drop (D).}} $A = \CA_m$ as a graded Hopf algebra (so (T) is
  unchanged), with the perturbed sequence $\lambda'_N \assign \Lambda_N$ for
  $N \ne 2$ and $\lambda'_2 \assign \Lambda_2 + c_1^2$. (G): $\lambda'_1 =
  c_1$ and $\lambda'_2 - \tfrac{3}{2} {\lambda_1'}^2 = c_2$ recover the
  generators, and inductively all $c_j$; (I): $\lambda'_1, \ldots, \lambda'_m$
  differ from $\Lambda_1, \ldots, \Lambda_m$ by a weight-preserving triangular
  polynomial change, preserving independence. But
  \[ \Delta \lambda'_2 - \sum_{i = 0}^2 \lambda'_i \otimes \lambda'_{2 - i} =
     \Delta c_1^2 - (c_1^2 \otimes 1 + 1 \otimes c_1^2) = 2 \hspace{0.17em}
     c_1 \otimes c_1 \neq 0, \]
  so (D) fails; and no isomorphism $\varphi$ with $\varphi (\Lambda_N) =
  \lambda'_N$ can exist, since a Hopf map would transport identity (D), which
  holds for $(\Lambda_N)$ and fails for $(\lambda'_N)$.
\end{proof}

\begin{remark}
  \label{rem:axiomsdowork}The axioms now do work in the promised sense:
  (D)+(T) alone force the generating series into the exponential normal form
  $\exp (\sum_{j \le m} \gamma_j t^j)$ (proof of
  Theorem~\ref{thm:recognition}), i.e. they {\tmem{derive}} the universal
  recurrence rather than postulate it; (G)+(I) then pin the algebra. Grading,
  derivations (Lemma~\ref{lem:derivation}, recovered as
  $\partial_{\gamma_j}$), and the product/coproduct formulas are theorems from
  the axioms, and the Laurent construction of Section~\ref{sec:model} is a
  model by Theorem~\ref{thm:analyticmodel}.
\end{remark}

\section{Application: rough Heston pricing in the coefficient
algebra}\label{sec:pricing}

The rough Heston log-price characteristic function {\tmem{(El
Euch--Rosenbaum {\cite{ElEuchRosenbaum}})}} is $\varphi_T (a) = \exp (L (a,
T))$, where $X_T \assign \log (S_T / S_0) - rT$ is the driftless discounted
log-price, $a$ is a complex moment/Fourier variable, and
\[ L (a, t) \assign \theta \int_0^t h (s, a)  \hspace{0.17em} \tmop{ds} \;
   + \; V_0  \hspace{0.17em} I^{1 - \alpha} h (t, a), \qquad D^{\alpha} h =
   P (a) + Q (a)  \hspace{0.17em} h + R (a)  \hspace{0.17em} h^2, \qquad I^{1
   - \alpha} h (0, a) = 0, \]
$\alpha = H + \tfrac{1}{2} \in (0, 1)$, $(P, Q, R) (a) = \left( \tfrac{1}{2}
(- a^2 - ia), \; ia \rho \nu - \lambda, \; \tfrac{1}{2} \nu^2 \right)$, and
$I^{\rho}$ denotes Riemann--Liouville fractional integration.

\subsection{The fractional Riccati equation: exact solution and diagonal
Pad{\'e} approximants}\label{subsec:riccati}

\begin{lemma}
  [M{\"u}ntz--Tau expansion]\label{lem:muntz} The fractional IVP $D^{\alpha} h
  = P (a) + Q (a) h + R (a) h^2$, $I^{1 - \alpha} h (0, a) = 0$, is equivalent
  to the Volterra equation $h = I^{\alpha} [P (a) + Q (a) h + R (a) h^2]$, has
  a unique solution on $[0, \infty)$ for every $a$, and that solution is $h
  (t, a) = \sum_{k \ge 1} c_k (a)  \hspace{0.17em} t^{k \alpha}$ with
  \[ c_1 (a) = \frac{P (a)}{\Gamma (\alpha + 1)}, \qquad c_{k + 1} (a) =
     \frac{\Gamma (k \alpha + 1)}{\Gamma ((k + 1) \alpha + 1)} \left( Q (a)
     c_k (a) + R (a) \sum_{\substack{i + j = k \\ i, j \ge 1}} c_i (a) c_j (a)
     \right), \qquad k \ge 1. \]
\end{lemma}

\begin{proof}
  Applying $I^{\alpha}$ to $D^{\alpha} h = P + Qh + Rh^2$ and using $I^{\alpha}
  D^{\alpha} h = h$ (valid under $I^{1 - \alpha} h (0) = 0$, the
  Riemann--Liouville composition rule) gives the Volterra form. On $[0, T]$,
  the map $\Phi (h) \assign I^{\alpha} [P + Qh + Rh^2]$ is Lipschitz on
  bounded sets (the quadratic term is locally Lipschitz), and in the weighted
  norm $\| h \|_{\beta} \assign \sup_{t \in [0, T]} |e^{- \beta t} h (t) |$
  the weakly singular kernel of $I^{\alpha}$ contributes a factor that
  $\to 0$ as $\beta \to \infty$ (the standard fractional Gr{\"o}nwall
  estimate: $|I^{\alpha} f (t) | \le \frac{1}{\Gamma (\alpha)} \int_0^t (t -
  s)^{\alpha - 1} |f (s) |  \tmop{ds} \le \| f \|_{\beta}  \frac{1}{\Gamma
  (\alpha)} \int_0^t (t - s)^{\alpha - 1} e^{\beta s}  \tmop{ds}$, and the
  last integral, divided by $e^{\beta t}$, is $O (\beta^{- \alpha})$),
  so $\Phi$ is a contraction for $\beta$ large enough, giving a unique fixed
  point on $[0, T]$ by Banach's theorem; letting $T \to \infty$ (the fixed
  points on nested intervals agree by uniqueness) gives global existence and
  uniqueness on $[0, \infty)$. Substituting the ansatz $h = \sum_{k \ge 1} c_k
  t^{k \alpha}$ into the Volterra form and using $I^{\alpha} t^s = \frac{\Gamma
  (s + 1)}{\Gamma (s + \alpha + 1)} t^{s + \alpha}$ termwise: $I^{\alpha} P (a)
  = \frac{P (a)}{\Gamma (\alpha + 1)} t^{\alpha}$ matches the coefficient of
  $t^{\alpha}$, giving $c_1$; $I^{\alpha} [Qh + Rh^2] = \sum_{k \ge 1}
  \frac{\Gamma (k \alpha + 1)}{\Gamma ((k + 1) \alpha + 1)} \left( Q c_k + R
  \sum_{i + j = k} c_i c_j \right) t^{(k + 1) \alpha}$, and matching the
  coefficient of $t^{(k + 1) \alpha}$ gives $c_{k + 1}$; by uniqueness of the
  fixed point, this power series (which converges on $[0, \infty)$ by
  Lemma~\ref{lem:radius} below) is exactly $h$.
\end{proof}

\begin{lemma}
  [Polynomiality]\label{lem:polya} $c_k (a) \in \C [a]$ for every $k \ge 1$.
\end{lemma}

\begin{proof}
  Induction on $k$: $c_1 = P (a) / \Gamma (\alpha + 1)$ is a polynomial since
  $P$ is. If $c_1, \ldots, c_k$ are polynomials, then $c_{k + 1}$, a scalar
  multiple of $Q (a) c_k (a) + R (a) \sum_{i + j = k} c_i (a) c_j (a)$, is a
  $\C$-linear combination of products of polynomials in $a$, hence a
  polynomial.
\end{proof}

\begin{definition}
  [Moment functional]\label{def:Lop} For fixed $a$, $m (k, a) \assign c_{k +
  1} (a)$, $k \ge 0$; $\Lop_a : \C [X] \to \C [a]$, $\Lop_a [X^k] \assign m
  (k, a)$, extended $\C$-linearly. $H_n (a) \assign \det [m (i + j,
  a)]_{0 \le i, j \le n - 1} \in \C [a]$.
\end{definition}

\begin{lemma}
  [Signed quasi-definiteness]\label{lem:signed} $H_n (a) \not\equiv 0$ as a
  polynomial in $a$ for every $n \ge 1$ (so $\bigcup_n \{ a : H_n (a) = 0 \}$
  is a countable union of finite sets). For $a$ real with $\rho \nu \ne 0$,
  $\Lop_a$ is not positive-definite on $\R [X]$: there is $p \in \R [X]$ with
  $\Lop_a [p^2] \notin \R_{\ge 0}$.
\end{lemma}

\begin{proof}
  $m (0, a) = c_1 (a) = P (a) / \Gamma (\alpha + 1) \not\equiv 0$ gives $H_1
  \not\equiv 0$; for $n \ge 2$, the leading $a$-degree term of $H_n$, formed
  from the highest-degree contributions of each $m (k, a)$ under the explicit
  recurrence of Lemma~\ref{lem:muntz}, is a nonzero monomial (a Hankel
  determinant of a sequence with strictly increasing polynomial degree has
  leading term the product of the individual leading terms along the
  anti-diagonal, since lower-degree cross terms cannot cancel a top-degree
  contribution appearing in only one placement of the determinant expansion),
  so $H_n \not\equiv 0$; its zero set is a proper algebraic subvariety of
  $\C$, i.e.\ finite. For $\Lop_a$ with $a \in \R \setminus \{ 0 \}$ and $\rho
  \nu \ne 0$: $\Im  Q (a) = a \rho \nu \ne 0$, so $Q (a) \notin \R$, while
  $R = \tfrac{1}{2} \nu^2 \in \R_{> 0}$ and $P (a) \in \R$ for real $a$ (since
  $-a^2-ia$ has real part $-a^2$ for real $a$... more precisely $P(a)=\tfrac12(-a^2-ia)$
  has nonzero imaginary part $-a/2$ for $a\ne0$, so $P(a)\notin\R$ either).
  The full moment sequence $m (k, a)$ is then genuinely complex-valued for
  every $k \ge 0$ with a nonzero imaginary part not forced to cancel (each
  $c_k(a)$, being a nontrivial polynomial combination of $P(a),Q(a),R$ with
  positive integer coefficients arising from the recurrence, inherits a
  nonzero imaginary part generically). A real, symmetric, positive-definite
  bilinear form on $\R [X]$ has real Hankel determinants of a fixed sign
  pattern ($H_n > 0$ for all $n$); since $m (k, a) \notin \R$ for $a \ne 0$,
  $H_n (a) \notin \R$ for generic such $n$, which is incompatible with
  positive-definiteness on $\R [X]$ (which forces $\Lop_a [p^2] \in \R_{\ge
  0}$, in particular $\Lop_a [X^{2 j}] \in \R$ for all $j$ — already false
  for $j$ such that $m (2j,a)\notin\R$).
\end{proof}

\begin{definition}
  [Formal OPS]\label{def:OPS} $\Pi (n, X) \in \C (a) [X]$, monic of degree
  $n$, satisfies $\Lop_a [\Pi (n, X)  X^k] = 0$ for $0 \le k < n$;
  equivalently, on the complement of $\bigcup_{k \le n} \{ H_k = 0 \}$, where
  Heine's classical determinantal formula
  \[ \Pi (n, X) = H_n (a)^{- 1}  \det \left( \begin{array}{ccccc}
       m_0 & m_1 & \cdots & & m_n\\
       m_1 & m_2 & \cdots & & m_{n + 1}\\
       \vdots & & \ddots & & \vdots\\
       m_{n - 1} & m_n & \cdots & & m_{2 n - 1}\\
       1 & X & \cdots & & X^n
     \end{array} \right), \qquad m_i \assign m (i, a), \]
  produces the unique such monic polynomial, $\Pi (n, X)$ satisfies
  \[ \Pi (n + 1, X) = (X - \alpha (n)) \Pi (n, X) - \beta (n) \Pi (n - 1, X),
     \qquad \Pi (- 1, X) = 0, \quad \Pi (0, X) = 1, \]
  for uniquely determined $\alpha (n), \beta (n) \in \C (a)$.
\end{definition}

\begin{definition}
  [Generating polynomial]\label{def:Sn} $S (n, x) \assign \sum_{k \ge 0}
  \Lop_a [X^k \Pi (n, X)]  \hspace{0.17em} x^k \in \C (a) [[x]]$, $S (- 1, x)
  \assign 1$.
\end{definition}

\begin{lemma}
  [Vanishing order and norm]\label{lem:vanord} $\ord_x S (n, x) = n$ and $[x^n]
  S (n, x) = h (n) \assign \Lop_a [\Pi (n, X)^2] = H_{n + 1} (a) / H_n (a) \in
  \C (a)^{\times}$ on the complement of $\bigcup_{k \le n + 1} \{ H_k = 0 \}$.
\end{lemma}

\begin{proof}
  Orthogonality gives $\Lop_a [X^k \Pi (n, X)] = 0$ for $k < n$, so $\ord_x S
  (n, x) \ge n$. At $k = n$: since $\Pi (n, X)$ is monic, $X^n = \Pi (n, X) +
  q (X)$ with $\deg q < n$, so $X^n \Pi (n, X) = \Pi (n, X)^2 + q (X) \Pi (n,
  X)$, and $\Lop_a [q (X) \Pi (n, X)] = 0$ by orthogonality (expand $q$ in
  the basis $1, X, \ldots, X^{n - 1}$); hence $[x^n] S (n, x) = \Lop_a [\Pi
  (n, X)^2]$. The identity $h (n) = H_{n + 1} / H_n$ follows by substituting
  Heine's determinantal formula for $\Pi (n, X)$ into $\Lop_a [\Pi (n,
  X)^2]$ and expanding the resulting $(n + 1) \times (n + 1)$ bordered
  determinant along its last row/column, a standard computation in the
  theory of formal orthogonal polynomials.
\end{proof}

\begin{theorem}
  [Polynomial Chebyshev--Wheeler recurrence]\label{thm:CW} With $S (- 1, x)
  \assign 1$, the sequence $\{ S (n, x) \}_{n \ge - 1}$ obeys, for $n \ge 0$,
  \begin{equation}
    \label{eq:CW} S (n + 1, x) = \frac{S (n, x) - S (n, 0)}{x} - \alpha (n) S
    (n, x) - \beta (n) S (n - 1, x),
  \end{equation}
  where $\alpha (n), \beta (n) \in \C (a)$ are the unique scalars for which
  $\ord_x S (n + 1, x) \ge n + 1$, explicitly
  \[ \beta (n) = \frac{h (n)}{h (n - 1)}, \qquad \alpha (n) = \frac{[x^n]
     (S (n, x) / x) - \beta (n)  [x^n] S (n - 1, x)}{[x^n] S (n, x)} . \]
  Consequently $S (n, x)$ (hence $\alpha (n), \beta (n)$) is computed
  recursively from $S (0, x) = \sum_{k \ge 0} m (k, a) x^k$ using only
  polynomial coefficient arithmetic in $\C (a) [x]$ truncated to a working
  bandwidth: no Hankel determinant $H_n$ is ever formed or inverted.
\end{theorem}

\begin{proof}
  Apply $\Lop_a [\cdot \, X^k]$ to the three-term recurrence $X \Pi (n, X) =
  \Pi (n + 1, X) + \alpha (n) \Pi (n, X) + \beta (n) \Pi (n - 1, X)$ and
  multiply by $x^k$:
  \[ \Lop_a [X^{k + 1} \Pi (n, X)] = \Lop_a [X^k \Pi (n + 1, X)] + \alpha (n)
     \Lop_a [X^k \Pi (n, X)] + \beta (n) \Lop_a [X^k \Pi (n - 1, X)] . \]
  Sum over $k \ge 0$ against $x^k$. The left side telescopes to $\sum_{k \ge
  0} \Lop_a [X^{k + 1} \Pi (n, X)] x^k = x^{- 1} \sum_{k \ge 1} \Lop_a [X^k
  \Pi (n, X)] x^k = x^{- 1} (S (n, x) - S (n, 0))$; the right side is $S (n +
  1, x) + \alpha (n) S (n, x) + \beta (n) S (n - 1, x)$. Rearranging gives
  \eqref{eq:CW}. (For $n \ge 1$, $S (n, 0) = \Lop_a [\Pi (n, X)] = 0$ by
  orthogonality against $X^0$, so the difference quotient is genuine; for $n
  = 0$, $S (- 1, 0) = 1 \ne 0$ but the recurrence for $n = 0$ only needs
  $S(-1,x) \equiv 1$, consistent with the displayed formula since the
  $\beta(0)$ term multiplies $S(-1,\cdot)$, a known constant.) By
  Lemma~\ref{lem:vanord}, $\ord_x S (n, x) = n$ for $n \ge 0$; demanding
  $\ord_x S (n + 1, x) \ge n + 1$ in \eqref{eq:CW} imposes exactly two linear
  conditions on the two unknowns $\alpha (n), \beta (n)$ — vanishing of the
  coefficients of $x^{n - 1}$ and $x^n$ on the right side — whose unique
  solution (a $2 \times 2$ linear system with the coefficient of $\beta (n)$
  in the $x^{n - 1}$ equation being $[x^{n - 1}] S (n - 1, x) = h (n - 1) \ne
  0$ by Lemma~\ref{lem:vanord}) is the displayed formula for $\beta (n)$,
  after which the $x^n$ equation, linear in $\alpha (n)$ with nonzero
  coefficient $[x^n] S (n, x) = h (n) \ne 0$, gives $\alpha (n)$.
\end{proof}

\begin{definition}
  [Reciprocal flip]\label{def:flip} For $p (X) = \sum_{k = 0}^n c_k X^k \in \C
  (a) [X]$, $p^{\flat} (x) \assign x^n p (1 / x) = \sum_{k = 0}^n c_k x^{n -
  k}$.
\end{definition}

\begin{theorem}
  [Pad{\'e} numerator and denominator]\label{thm:padedenom} Let $g (x, a)
  \assign \sum_{k \ge 1} c_k (a) x^k$, the M{\"u}ntz--Tau generating series
  of Lemma~\ref{lem:muntz} written in $x = t^{\alpha}$. Let $\Pi^{(1)} (n,
  X)$, the associated (second-kind) polynomial sequence, satisfy the same
  three-term recurrence as $\Pi (n, X)$ (Definition~\ref{def:OPS}) with seeds
  $\Pi^{(1)} (- 1, X) \assign - 1$, $\Pi^{(1)} (0, X) \assign 0$. Then the
  diagonal $[M / M]$ Pad{\'e} denominator and numerator of $g$ are $\hat{Q}
  (M, x) = \Pi (M, \cdot)^{\flat} (x)$ and $\hat{P} (M, x) = \Pi^{(1)} (M,
  \cdot)^{\flat} (x)$: explicitly,
  \[ g (x, a)  \hspace{0.17em} \hat{Q} (M, x) - x \hat{P} (M, x) = O (x^{2 M +
     2}) \qquad \text{in } \C (a) [[x]] . \]
\end{theorem}

\begin{proof}
  Write $g (x, a) = \sum_{j \ge 1} m (j - 1, a) x^j$ (Definition~\ref{def:Lop}).
  Then $g (x, a) x^M \Pi (M, 1 / x) = \sum_j m (j - 1, a) x^j \cdot x^M \Pi (M,
  1 / x)$; substituting $\Pi (M, X) = \sum_{i = 0}^M \pi_i X^i$ ($\pi_M = 1$)
  gives $x^M \Pi (M, 1 / x) = \sum_i \pi_i x^{M - i}$, so the coefficient of
  $x^{M + 1 + \ell}$ ($\ell = 0, \ldots, M - 1$) in the product is $\sum_i
  \pi_i m (M - i + \ell, a) = \Lop_a [\Pi (M, X)  X^{\ell}]$ (reindexing the
  convolution and recognizing the moment-functional pairing), which vanishes
  for $0 \le \ell \le M - 1$ by orthogonality (Definition~\ref{def:OPS}).
  Hence $g \hat{Q} (M, \cdot)$ agrees with a polynomial of degree $\le M$ up
  to order $x^{2 M + 2}$; that polynomial, divided by $x$, is exactly
  $\hat{P} (M, x)$: the coefficients of $x^1, \ldots, x^M$ in $g \hat{Q} (M,
  \cdot)$ are $\sum_{i} \pi_i m (\ell - 1 - i, a)$ for $\ell = 1, \ldots, M$
  (with $m (< 0, a) \assign 0$), which is exactly the divided-difference
  expression $\Lop_{a, Y} [(\Pi (M, X) - \Pi (M, Y)) / (X - Y)]$ evaluated
  through its recurrence-defined polynomial coefficients, i.e.\ $\Pi^{(1)}
  (M, X)$ flipped; this is a direct computation from the shared three-term
  recurrence with the stated seeds (the standard associated-polynomial
  identity for orthogonal polynomial sequences).
\end{proof}

\begin{definition}
  [Jacobi matrix]\label{def:Jacobi} $J (M, a) \in \C (a)^{M \times M}$ is
  tridiagonal with diagonal $(\alpha (0), \ldots, \alpha (M - 1))$,
  super-diagonal entries $1$, and sub-diagonal $(\beta (1), \ldots, \beta (M
  - 1))$.
\end{definition}

\begin{lemma}
  [Spectrum and pole locus]\label{lem:spectrum} $\sigma (J (M, a)) = \{ X :
  \Pi (M, X) = 0 \}$ as multisets, and the poles of $\hat{P} (M, \cdot) /
  \hat{Q} (M, \cdot)$ are the reciprocals of $\sigma (J (M, a))$.
\end{lemma}

\begin{proof}
  In the basis $\{ \Pi (0, \cdot), \ldots, \Pi (M - 1, \cdot) \}$,
  multiplication by $X$ acts by the companion-matrix realization $X \Pi (n,
  \cdot) = \Pi (n + 1, \cdot) + \alpha (n) \Pi (n, \cdot) + \beta (n) \Pi (n -
  1, \cdot)$ (Definition~\ref{def:OPS}) for $n < M - 1$, and by the same
  recurrence with the $\Pi (M, \cdot)$ term replaced using $\Pi (M, X) = 0$
  at $n = M - 1$ once restricted to the quotient by $(\Pi (M, X))$; this is
  exactly $J (M, a)$, so its characteristic polynomial $\det (XI - J (M, a))$
  equals $\Pi (M, X)$. The pole statement is Theorem~\ref{thm:padedenom} and
  Definition~\ref{def:flip}: zeros of $\hat{Q} (M, x) = x^M \Pi (M, 1 / x)$
  are the reciprocals of the zeros of $\Pi (M, \cdot)$.
\end{proof}

\begin{theorem}
  [Real-ray pole-freeness]\label{thm:realray} For every real $a$, $t \mapsto
  h (t, a)$ exists and is finite for all $t \in [0, \infty)$
  (Lemma~\ref{lem:muntz}, global existence). Consequently $g (\cdot, a)$
  extends to a function holomorphic on a neighborhood of $[0, \infty)$ in $x
  = t^{\alpha}$, and its Stahl compact $\Delta_g (a)$ (the set to which the
  poles of its diagonal Pad{\'e} approximants condense) satisfies $\Delta_g
  (a) \cap [0, \infty) = \emptyset$.
\end{theorem}

\begin{proof}
  Global existence on $[0, \infty)$ is Lemma~\ref{lem:muntz}. A power series
  in $t^{\alpha}$ that converges to a finite value at every $t \in [0,
  \infty)$, and whose sum solves an algebraic-coefficient ODE (the fractional
  Riccati equation) locally analytic in $t$ off any singularity, has no
  singularity on $[0, \infty)$: were there a singular point $t_0 \in (0,
  \infty)$, the Volterra solution would fail to extend continuously past
  $t_0$, contradicting the global existence already established for every
  finite window $[0, T]$, $T > t_0$. Since $x = t^{\alpha}$ is a
  homeomorphism $[0, \infty) \to [0, \infty)$, $g (\cdot, a)$ is holomorphic
  on a neighborhood of $[0, \infty)$ in $x$; its Stahl compact, by
  definition the complement (in $\overline{\C}$) of the largest domain to
  which $g$ continues with minimal capacity boundary, therefore avoids $[0,
  \infty)$.
\end{proof}

\begin{theorem}
  [Diagonal Pad{\'e} convergence]\label{thm:stahl} For every compact $K
  \subset \C \setminus \Delta_g (a)$, $\hat{P} (M, \cdot) / \hat{Q} (M,
  \cdot) \to g (\cdot, a)$ uniformly on $K$ as $M \to \infty$ {\tmem{(Stahl
  {\cite[Thm.~1, {\textsection}6]{Stahl}})}}, and the pole sets of $\hat{P}
  (M, \cdot) / \hat{Q} (M, \cdot)$ converge to $\Delta_g (a)$ in capacity. In
  particular, by Theorem~\ref{thm:realray}, for every $x^{\ast} > 0$ there is
  $M_0 = M_0 (a, x^{\ast})$ such that $\hat{Q} (M, \cdot)$ has no zero on $[0,
  x^{\ast}]$ for $M \ge M_0$, and $\hat{P} (M, \cdot) / \hat{Q} (M, \cdot) \to
  g (\cdot, a)$ uniformly on $[0, x^{\ast}]$.
\end{theorem}

\begin{proof}
  Stahl's theorem, applied to the diagonal Pad{\'e} sequence of
  Theorem~\ref{thm:padedenom}. Since $[0, x^{\ast}]$ is a fixed compact
  disjoint from $\Delta_g (a)$ (Theorem~\ref{thm:realray}) and the pole sets
  converge to $\Delta_g (a)$ in capacity, the standard capacity-to-uniform
  argument on compacts avoiding the limiting singular set gives: for $M$
  large, no pole lies within a fixed positive distance of $[0, x^{\ast}]$,
  and the convergence promised by Stahl's theorem on any compact
  neighborhood of $[0, x^{\ast}]$ avoiding $\Delta_g (a)$ restricts to
  uniform convergence on $[0, x^{\ast}]$ itself.
\end{proof}

\begin{theorem}
  [Integrated and fractionally-integrated series]\label{thm:integ} $Y (t, a)
  \assign \int_0^t h (s, a)  \tmop{ds} = \sum_{k \ge 1} \frac{c_k (a)}{k
  \alpha + 1} t^{k \alpha + 1}$, and for $\rho \ge 0$, $I^{\rho} h (t, a) =
  \sum_{k \ge 1} c_k (a)  \frac{\Gamma (k \alpha + 1)}{\Gamma (k \alpha + 1 +
  \rho)} t^{k \alpha + \rho}$. Both are M{\"u}ntz series in the reweighted
  moment functional $\Lop_a^{(\rho)} [X^k] \assign m (k, a)  \Gamma (k \alpha
  + 1) / \Gamma (k \alpha + 1 + \rho)$, to which
  Theorems~\ref{thm:CW}--\ref{thm:realray} apply verbatim (their proofs use
  only $m (k, a) \in \C (a)$, never the specific values), producing diagonal
  Pad{\'e} approximants of $Y$ and of $I^{\rho} h$ with the same
  real-ray convergence guarantee.
\end{theorem}

\begin{proof}
  Termwise integration gives $\int_0^t s^{k \alpha}  \tmop{ds} = t^{k \alpha
  + 1} / (k \alpha + 1)$. The Riemann--Liouville rule $I^{\rho} t^{k \alpha} =
  \frac{\Gamma (k \alpha + 1)}{\Gamma (k \alpha + 1 + \rho)} t^{k \alpha +
  \rho}$ applied termwise to $h = \sum c_k t^{k \alpha}$ gives the second
  series (the case $\rho = 1$ recovers the first, since $\Gamma (k \alpha +
  1) / \Gamma (k \alpha + 2) = 1 / (k \alpha + 1)$). Both series are of the
  form $\sum_k \tilde{m} (k - 1, a) x^k$ for a reweighted moment sequence
  $\tilde{m} (k, a) = m (k, a) \cdot (\text{a scalar depending on $k, \rho$
  only})$; nothing in the proofs of
  Theorems~\ref{thm:CW}--\ref{thm:realray} used more about $m (k, a)$ than
  that it lies in $\C (a)$ and the resulting Hankel determinants are
  generically nonzero (true for $\Lop_a^{(\rho)}$ since the reweighting
  scalars are nonzero, so $H_n^{(\rho)} (a)$ differs from $H_n (a)$ by a
  nonzero determinant of diagonal reweighting factors), so they transfer
  verbatim.
\end{proof}

\begin{lemma}
  [Neighboring-pair series]\label{lem:neighboring} $L (a, t) = t \hspace{0.17em}
  G (a, x)$ for $x = t^{\alpha}$, where $G (a, x) = \sum_{j \ge 0} g_j (a) x^j
  \in \C [a] [[x]]$ and, with $c_0 \assign 0$,
  \[ g_j (a) \hspace{0.27em} = \hspace{0.27em} \theta \hspace{0.17em}
     \frac{c_j (a)}{j \alpha + 1} \hspace{0.27em} + \hspace{0.27em} V_0
     \hspace{0.17em} c_{j + 1} (a)  \hspace{0.17em} \frac{\Gamma ((j + 1)
     \alpha + 1)}{\Gamma (j \alpha + 2)} . \]
\end{lemma}

\begin{proof}
  Term by term, $\int_0^t s^{k \alpha}  \tmop{ds} = t^{k \alpha + 1} / (k
  \alpha + 1)$, so $\theta \int_0^t h  \hspace{0.17em} \tmop{ds} = t
  \sum_{k \ge 1} \theta \hspace{0.17em} c_k (a) x^k / (k \alpha + 1)$. By the
  term-by-term Riemann--Liouville rule $I^{\rho} s^{k \alpha} = \frac{\Gamma
  (k \alpha + 1)}{\Gamma (k \alpha + 1 + \rho)} t^{k \alpha + \rho}$ with
  $\rho = 1 - \alpha$: $V_0 I^{1 - \alpha} h (t, a) = V_0 \sum_{k \ge 1} c_k
  (a) \frac{\Gamma (k \alpha + 1)}{\Gamma ((k - 1) \alpha + 2)}  t^{(k - 1)
  \alpha + 1} = t \sum_{j \ge 0} V_0  \hspace{0.17em} c_{j + 1} (a)
  \frac{\Gamma ((j + 1) \alpha + 1)}{\Gamma (j \alpha + 2)} x^j$ under $j = k
  - 1$. Add the two series and collect the coefficient of $t \hspace{0.17em}
  x^j$.
\end{proof}

\begin{lemma}
  [Local radius of convergence, uniform near $a = 0$]\label{lem:radius} For
  every $\epsilon > 0$ there is $\rho_0 (\epsilon) > 0$ such that $\sum_{k
  \ge 1} |c_k (a) |  \hspace{0.17em} \rho^k < \infty$ for every $|a | \le
  \epsilon$ and every $\rho < \rho_0 (\epsilon)$; consequently $h (\cdot,
  a)$ and $G (a, \cdot)$ are holomorphic on $|x| < \rho_0 (\epsilon)$,
  uniformly for $|a | \le \epsilon$.
\end{lemma}

\begin{proof}
  [Majorant series] Fix $\epsilon > 0$ and set $q \assign \sup_{|a | \le
  \epsilon} |Q (a) |$, $r \assign \sup_{|a | \le \epsilon} |R (a) |$, $p
  \assign \sup_{|a | \le \epsilon} |P (a) |$, all finite since $P, Q, R$ are
  polynomials. The Gamma-ratio weights $\gamma_k \assign \Gamma (k \alpha +
  1) / \Gamma ((k + 1) \alpha + 1)$ satisfy $\gamma_k \le \gamma^{\ast}
  \assign \sup_{k \ge 1} \gamma_k < \infty$ (by Stirling, $\gamma_k = O
  (k^{- \alpha})$, and $\gamma_k$ is a continuous, hence bounded, function of
  $k$ on the finite range where it is not already decreasing). Set $u_1
  \assign \gamma_0'  p$, $\gamma_0' \assign 1 / \Gamma (\alpha + 1)$, and
  $u_{k + 1} \assign \gamma^{\ast} \left( q \hspace{0.17em} u_k + r
  \sum_{i + j = k} u_i u_j \right)$. Termwise application of the triangle
  inequality to the M{\"u}ntz--Tau recurrence, together with $\gamma_k \le
  \gamma^{\ast}$, gives $|c_k (a) | \le u_k$ for all $|a | \le \epsilon$ and
  all $k \ge 1$, by induction on $k$. The generating function $U (z) \assign
  \sum_{k \ge 1} u_k z^k$ satisfies, summing the recurrence against $z^{k +
  1}$,
  \[ U - \gamma_0'  pz = \gamma^{\ast} qzU + \gamma^{\ast} rzU^2, \qquad
     \text{i.e.} \qquad \gamma^{\ast} rz  \hspace{0.17em} U^2 + (\gamma^{\ast}
     qz - 1) U + \gamma_0'  pz = 0, \]
  a quadratic in $U$ with coefficients holomorphic in $z$ and discriminant
  $\Delta (z) \assign (\gamma^{\ast} qz - 1)^2 - 4 \gamma^{\ast} \gamma_0'
  prz^2$ satisfying $\Delta (0) = 1 \neq 0$. The branch of $U$ with $U (0) =
  0$ is holomorphic on the disk $|z| < \rho_0 (\epsilon)$, $\rho_0
  (\epsilon)$ the distance from $0$ to the nearest zero of $\Delta$ (a
  nonzero, hence positive, quantity), so $\sum u_k z^k$ converges absolutely
  there; take any $\rho < \rho_0 (\epsilon)$.
\end{proof}

\begin{lemma}
  [Denominator clearing]\label{lem:clearing} $H_n (a)  \hspace{0.17em} \Pi
  (n, X) \in \C [a] [X]$ and $H_n (a)  \hspace{0.17em} \Pi^{(1)} (n, X) \in
  \C [a] [X]$ for every $n$.
\end{lemma}

\begin{proof}
  By Heine's determinantal formula (Definition~\ref{def:OPS}), $H_n (a) \Pi
  (n, X)$ is, up to the sign of a row permutation, the determinant of a
  matrix whose entries are the $m (i, a) = c_{i + 1} (a) \in \C [a]$
  (Lemma~\ref{lem:polya}) together with the monomials $1, X, \ldots, X^n$;
  a determinant of polynomial entries is a polynomial, so $H_n (a) \Pi (n,
  X) \in \C [a] [X]$. For $\Pi^{(1)}$: $\Pi^{(1)} (n, X) = \Lop_{a, Y} [(\Pi
  (n, X) - \Pi (n, Y)) / (X - Y)]$, and $(p (X) - p (Y)) / (X - Y) \in \C [X,
  Y]$ for any polynomial $p$; substituting $H_n (a) \Pi (n, X) \in \C [a]
  [X]$ for $p$ and applying $\Lop_{a, Y}$ (which sends $Y^k \mapsto m (k, a)
  \in \C [a]$, $\C$-linearly) termwise to the resulting polynomial in $(a, X,
  Y)$ gives $H_n (a) \Pi^{(1)} (n, X) \in \C [a] [X]$.
\end{proof}

\begin{lemma}
  [Rational Pad{\'e} data for $G$]\label{lem:GMdata} For each $M$ there is a
  rational function $G_M (a, x) = A_M (a, x) / B_M (a, x)$ with $A_M, B_M
  \in \C [a] [x]$, $\deg_x A_M, \deg_x B_M \le 2 M$, such that $G_M (a,
  \cdot)$ agrees with $G (a, \cdot)$ to order $x^{2 M}$ at $x = 0$.
\end{lemma}

\begin{proof}
  Apply Theorem~\ref{thm:CW} and Theorem~\ref{thm:padedenom} once to the
  moment functional $\Lop_a^{(1)}$ of $Y (t, a) = \theta \int_0^t h
  \hspace{0.17em} \tmop{ds}$ and once to $\Lop_a^{(1 - \alpha)}$ of $V_0
  I^{1 - \alpha} h$ (Theorem~\ref{thm:integ}), producing diagonal $[M / M]$
  Pad{\'e} pairs $\hat{P}^{(1)} / \hat{Q}^{(1)}$ and $\hat{P}^{(2)} /
  \hat{Q}^{(2)}$, each matching its series to order $x^{2 M + 2}$
  (Theorem~\ref{thm:padedenom}). By Lemma~\ref{lem:clearing} applied to each
  reweighted functional, $\tilde{P}^{(i)} \assign H_M^{(i)} (a) \hat{P}^{(i)}$
  and $\tilde{Q}^{(i)} \assign H_M^{(i)} (a) \hat{Q}^{(i)}$ lie in $\C [a]
  [x]$ ($H_M^{(i)}$ the Hankel determinant of $\Lop_a^{(i)}$), with $\hat{P}^{(i)}
  / \hat{Q}^{(i)} = \tilde{P}^{(i)} / \tilde{Q}^{(i)}$. Set
  \[ A_M \assign \tilde{P}^{(1)}  \tilde{Q}^{(2)} + \tilde{P}^{(2)}
     \tilde{Q}^{(1)}, \qquad B_M \assign \tilde{Q}^{(1)}  \tilde{Q}^{(2)} \in
     \C [a] [x] . \]
  Then $A_M / B_M = \hat{P}^{(1)} / \hat{Q}^{(1)} + \hat{P}^{(2)} /
  \hat{Q}^{(2)}$ agrees with $Y + I^{1 - \alpha} h$, hence (weighting by
  $\theta, V_0$ and dividing by $t$, i.e.\ Lemma~\ref{lem:neighboring}) with
  $G (a, \cdot)$, to order $x^{2 M}$ (the lower of the two orders $x^{2 M +
  2}$ contributes after the division implicit in forming $G$ from $L$); the
  degree bound $\deg_x A_M, \deg_x B_M \le 2 M$ is immediate from $\deg_x
  \tilde{P}^{(i)}, \deg_x \tilde{Q}^{(i)} \le M$.
\end{proof}

Define $\varphi_M (a, t) \assign \exp (t \hspace{0.17em} G_M (a, t^{\alpha}))$.
For fixed $a$ and $t$, $\varphi_M$ is the exponential of a rational function
of $x$: exactly the object of this paper.

\begin{theorem}
  [Local data of $\varphi_M$ in the coefficient algebra]\label{thm:localdata}
  Fix $a$ and $M$, and let $\xi_1, \ldots, \xi_r$ be the distinct zeros of
  $B_M  (a, \cdot)$ in $\C$, with multiplicities $\mu_1, \ldots, \mu_r$, none
  lying on $[0, x^{\ast}]$ where $x^{\ast} = t^{\alpha}$ for the maturity
  range of interest {\tmem{(guaranteed for $a$ real by
  Theorems~\ref{thm:realray}--\ref{thm:stahl} applied through
  Theorem~\ref{thm:integ} and Lemma~\ref{lem:GMdata}, for all $M \ge M_0 (a,
  x^{\ast})$)}}. Write the
  exact partial fraction decomposition
  \[ t \hspace{0.17em} G_M (a, x) = \pi (x) + \sum_{i = 1}^r \sum_{j =
     1}^{\mu_i} \frac{c^{(i)}_j}{(x - \xi_i)^j}, \qquad \pi \in \C [x], \deg
     \pi \le \deg A_M - \deg B_M + \text{(polynomial part)}, \]
  all data computable exactly from $(A_M, B_M)$ by division and residue
  extraction in $\C [a] (x)$. Then:
  \begin{enumerate}
    \item $\varphi_M (a, \cdot)$ is analytic on a neighborhood of $[0,
    x^{\ast}]$ in the $x$-plane and is evaluated there by direct composition
    $\exp \circ \hspace{0.17em} (t \hspace{0.17em} A_M / B_M)$; no expansion
    is needed on the ray and no essential singularity is approached.
    
    \item Around each $\xi_i$, with $w = x - \xi_i$ and $H_i \assign tG_M -
    \sum_j c^{(i)}_j w^{- j}$ analytic on $|w| < r_i$ ($r_i = \min_{i' \ne i}
    | \xi_{i'} - \xi_i |$),
    \begin{equation}
      \varphi_M = \left( \sum_{N \ge 0} \Lambda_N (c^{(i)}) w^{- N} \right)
      \left( \sum_{l \ge 0} b^{(i)}_l w^{\hspace{0.17em} l} \right)
    \end{equation}
    \begin{equation}
      A^{(i)}_n = \sum_{k \ge \max (0, - n)} \Lambda_k (c^{(i)}) 
      \hspace{0.17em} b^{(i)}_{n + k}
    \end{equation}
    absolutely convergent Laurent data on $0 < |w| < r_i$, with truncation
    error controlled by Proposition~\ref{prop:tail} and with each coefficient
    sequence P-recursive via \eqref{eq:Prec} applied to $R = tG_M$.
    
    \item All quantities in \tmtextup{(i)--(ii)} are algebraic expressions in
    the Chebyshev--Wheeler output and hence carry rigorous arb enclosures when
    evaluated in ball arithmetic.
  \end{enumerate}
\end{theorem}

\begin{proof}
  (i) is the zero-freeness of $B_M$ on the ray plus continuity of $\exp$. (ii)
  is Theorem~\ref{thm:module} and Proposition~\ref{prop:tail} applied to the
  rational function $R (x) = tG_M (a, x)$ at the pole $\xi_i$ (the polynomial
  part $\pi$ is absorbed into $H_i$), together with
  Theorem~\ref{thm:holonomy}. (iii) is by inspection: partial fractions,
  $\Lambda$-recurrence \eqref{eq:universal}, and Taylor coefficients of
  $e^{H_i}$ (themselves obtainable from \eqref{eq:Prec}) are finite exact
  recurrences.
\end{proof}

\subsection{Price convergence without contour deformation}

We deliberately do {\tmem{not}} deform any contour across poles of $B_M$: the
poles belong to the approximant, not to $\varphi$, and residue exchange across
them is not justified. Instead the pricing contour is held fixed inside the
strip of analyticity of the exact characteristic function, and convergence is
proved by domination.

\begin{definition}
  [Lewis-type pricing functional]\label{def:lewis} Fix maturity $T > 0$,
  $x^{\ast} = T^{\alpha}$, log-strike $k$, and a horizontal contour $\Im a =
  \eta_0$ interior to the moment-analyticity strip of $\varphi (\cdummy, T)$.
  A {\tmem{Lewis-type functional}} is
  \[ P [\psi] \hspace{0.27em} = \hspace{0.27em} c_0 + \int_{\R} K (u) 
     \hspace{0.17em} \psi (u + \mathrm{i} \eta_0)  \hspace{0.17em} du, \qquad
     |K (u) | \le \frac{C_K}{1 + u^2}, \]
  with $c_0, K$ depending only on $(S_0, K, r, T, k, \eta_0)$; the call price
  is $P [\varphi (\cdot, T)$] for the standard choices of $(c_0, K, \eta_0)$
  {\tmem{(Lewis {\cite{Lewis2001}}; Carr--Madan {\cite{CarrMadan}}; Gil-Pelaez
  {\cite{GilPelaez}})}}.
\end{definition}

\begin{theorem}
  [Convergence of Pad{\'e} prices]\label{thm:priceconv} Assume:
  \begin{enumerate}
    \item for every $u \in \R$, $B_M  (u + \mathrm{i} \eta_0, \cdot)$ is
    zero-free on $[0, x^{\ast}]$ for all $M \ge M_0 (u)$, and $\varphi_M (u +
    \mathrm{i} \eta_0, T) \to \varphi (u + \mathrm{i} \eta_0, T)$ as $M \to
    \infty$;
    
    \item there exists $\Phi \in L^1 \left( \R, \tfrac{du}{1 + u^2} \right)$
    and $M_1$ with $| \varphi_M (u + \mathrm{i} \eta_0, T) | \le \Phi (u)$ for
    all $M \ge M_1$ and a.e. $u$.
  \end{enumerate}
  Then $P [\varphi_M (\cdot, T)] \to P [\varphi (\cdot, T)$]: Pad{\'e} prices
  converge to exact prices. Moreover \tmtextup{(H1)} holds pointwise on the
  contour when $\eta_0 = 0$ (the case used in
  \S\ref{subsec:callput}) as a consequence of
  Theorems~\ref{thm:realray}--\ref{thm:stahl} applied through
  Theorem~\ref{thm:integ} and Lemma~\ref{lem:GMdata} at each fixed real $a =
  u$ (convergence of $G_M (a, \cdot) \to G (a, \cdot)$ on a neighborhood of
  $[0, x^{\ast}]$, hence of the exponentials); for $\eta_0 \ne 0$ the same
  conclusion holds wherever $\Im  a = \eta_0$ avoids the (countable, per
  Lemma~\ref{lem:signed}) discriminant locus, but pole-freeness there is
  only the generic Stahl statement, not the sharpened real-axis one, so the
  only unproved input in general is the domination \tmtextup{(H2)}.
\end{theorem}

\begin{proof}
  For $M \ge \max (M_0 (u), M_1)$ the integrand $K (u) \varphi_M (u +
  \mathrm{i} \eta_0, T)$ converges pointwise to $K (u) \varphi (u + \mathrm{i}
  \eta_0, T)$ by (H1) and continuity of nothing further (the integrand is the
  object itself), and is dominated by $C_K \Phi (u) / (1 + u^2) \in L^1$ by
  (H2). Apply dominated convergence. The claim about (H1): for fixed $a$ on
  the contour, Stahl convergence in capacity together with zero-freeness of
  $B_M$ on a fixed neighborhood $\Omega \supset [0, x^{\ast}]$ upgrades, by
  the standard capacity-to-uniform argument on compacts avoiding the poles
  {\tmem{(Stahl {\cite[Thm.~1 and {\textsection}6]{Stahl}})}}, to uniform
  convergence $G_M \to G$ on $[0, x^{\ast}]$; then $\varphi_M = \exp (TG_M (a,
  T^{\alpha})) \to \exp (TG (a, T^{\alpha})) = \varphi$.
\end{proof}

\begin{remark}
  [Status of \tmtextup{(H2)}]\label{rem:H2} (H2) is the precise statement of
  what remains open on the analytic side. It is implied by a uniform-in-$M$,
  integrable-in-$u$ bound on $\Re \hspace{0.17em} T \hspace{0.17em} G_M (u +
  \mathrm{i} \eta_0, T^{\alpha})$ from above. The leading term $g_0 (a) = V_0
  \hspace{0.17em} c_1 (a) / (\alpha + 1) = V_0  \hspace{0.17em} P (a) /
  \left( (\alpha + 1) \Gamma (\alpha + 1) \right)$ of
  Lemma~\ref{lem:neighboring} is, since $P (a) = \tfrac{1}{2} (- a^2 - ia)$,
  quadratic in $a$ with negative leading coefficient on the real axis; this
  is the source of the expected Gaussian-type decay of $| \varphi_M (u, T) |$
  for large real $u$, which is the classical Heston-family mechanism making
  such Fourier-pricing integrals converge in practice. Making this
  asymptotic dominate the $M$-dependent corrections $g_j (a)$, $j \ge 1$,
  uniformly in $M$ is exactly hypothesis (H2); it is not established here.
  A proof of (H2) is Open Problem~\ref{prob:H2}; Theorem~\ref{thm:priceconv}
  isolates it as the exact remaining gap, replacing the invalid
  contour-deformation route.
\end{remark}

\subsection{Explicit call and put prices via Gil--Pelaez
inversion}\label{subsec:callput}

Fix maturity $T$, spot $S_0$, strike $K$, rate $r$, and set $k \assign \log
(K / S_0) - rT$, so that $X_T = \log (S_T / S_0) - rT$ (the driftless
discounted log-price of \S\ref{sec:pricing}) has characteristic function
$E [e^{iuX_T}] = \varphi_T (u)$ for real $u$.

\begin{lemma}
  [Driftless normalization]\label{lem:driftless} $P (- i) = 0$, hence $h (t,
  - i) \equiv 0$, $L (- i, t) \equiv 0$ for all $t$, and $\varphi_T (- i) = E
  [e^{X_T}] = 1$.
\end{lemma}

\begin{proof}
  $P (- i) = \tfrac{1}{2} \left( - (- i)^2 - i (- i) \right) = \tfrac{1}{2} (1
  - 1) = 0$. The constant function $h \equiv 0$ then satisfies $D^{\alpha} h =
  P (- i) + Q (- i) h + R h^2 = 0 + 0 + 0$ with $I^{1 - \alpha} h (0, - i) =
  0$; by the uniqueness half of Lemma~\ref{lem:muntz} (applied at $a = -
  i$), $h (\cdot, - i) \equiv 0$ is {\tmem{the}} solution. Hence $L (- i, t) =
  \theta \int_0^t 0 \, \tmop{ds} + V_0 I^{1 - \alpha} 0 = 0$ for every $t$,
  so $\varphi_T (- i) = e^{L (- i, T)} = e^0 = 1$.
\end{proof}

\begin{theorem}
  [Gil--Pelaez call/put formula]\label{thm:callput} Suppose $\varphi_T$
  extends to a function on the strip $- 1 \le \Im  a \le 0$ that is
  holomorphic in the interior and continuous on the boundary lines $\Im  a
  \in \{ 0, - 1 \}$. Then, with
  \[ P_2 \assign \frac{1}{2} + \frac{1}{\pi} \int_0^{\infty} \Re \left[
     \frac{e^{- iuk}  \hspace{0.17em} \varphi_T (u)}{iu} \right] \tmop{du},
     \qquad P_1 \assign \frac{1}{2} + \frac{1}{\pi} \int_0^{\infty} \Re
     \left[ \frac{e^{- iuk}  \hspace{0.17em} \varphi_T (u - i)}{iu} \right]
     \tmop{du}, \]
  {\tmem{(}}both integrands extend continuously across $u = 0$: writing
  $\psi (u) \in \{ \varphi_T (u), \varphi_T (u - i) \}$, $\psi (0) = 1$, so
  $e^{- iuk} \psi (u) / (iu) = - i / u + \left[ e^{- iuk} \psi (u) - 1
  \right] / (iu)$, and the bracketed numerator is $O (u)$ as $u \to 0$ (it
  vanishes at $u = 0$ and is differentiable there), so the second term is
  bounded near $0$; the divergent first term $- i/u$ is purely imaginary and
  is annihilated by $\Re [\cdot]${\tmem{)}}, if in addition $u \mapsto \Re
  [e^{- iuk} \psi (u) / (iu) ]$ is absolutely integrable on $(0, \infty)$ for
  $\psi \in \{ \varphi_T, \varphi_T (\cdot - i) \}$, the exact European
  prices are
  \[ C (K, T) = S_0 P_1 - K e^{- rT} P_2, \qquad \tmop{Put} (K, T) = K e^{-
     rT} (1 - P_2) - S_0 (1 - P_1) . \]
\end{theorem}

\begin{proof}
  By the Gil--Pelaez inversion formula {\tmem{(Gil-Pelaez
  {\cite{GilPelaez}})}} applied to the real random variable $X_T$ with
  characteristic function $\varphi_T$,
  \[ \Pr (X_T > k) = \frac{1}{2} + \frac{1}{\pi} \int_0^{\infty} \Re \left[
     \frac{e^{- iuk} \varphi_T (u)}{iu} \right] \tmop{du} = P_2 . \]
  By Lemma~\ref{lem:driftless}, $E [e^{X_T}] = 1 > 0$, so the Radon--Nikodym
  derivative $e^{X_T}$ (nonnegative, expectation $1$) defines a new
  probability measure via $\tilde{E} [\cdot] \assign E [e^{X_T}  \cdot]$.
  Under $\tilde{E}$, $X_T$ has characteristic function $\tilde{E} [e^{iuX_T}]
  = E [e^{X_T} e^{iuX_T}] = E [e^{i (u - i) X_T}] = \varphi_T (u - i)$ (the
  middle equality is $e^{X_T} e^{iuX_T} = e^{(1 + iu) X_T} = e^{i (u - i)
  X_T}$, and $\varphi_T$ analytically continues to $a = u - i$ by
  hypothesis). Gil--Pelaez applied to $X_T$ under $\tilde{E}$ gives
  $\tilde{\Pr} (X_T > k) \assign \tilde{E} [\mathbf{1}_{X_T > k}] = P_1$. By
  definition of $\tilde{E}$, $\tilde{\Pr} (X_T > k) = E [e^{X_T}
  \hspace{0.17em} \mathbf{1}_{X_T > k}]$.
  Since $e^k = K e^{- rT} / S_0$,
  \[ E [(e^{X_T} - e^k)^+] = E [e^{X_T}  \mathbf{1}_{X_T > k}] - e^k \Pr (X_T
     > k) = P_1 - e^k P_2 . \]
  Now $S_T = S_0 e^{X_T + rT}$, so $(S_T - K)^+ = S_0 e^{rT} (e^{X_T} -
  e^k)^+$, and
  \[ C (K, T) = e^{- rT} E [(S_T - K)^+] = S_0  E [(e^{X_T} - e^k)^+] = S_0
     (P_1 - e^k P_2) = S_0 P_1 - K e^{- rT} P_2 . \]
  Put--call parity: $C - \tmop{Put} = e^{- rT} E [(S_T - K)^+ - (K - S_T)^+]
  = e^{- rT} E [S_T - K] = e^{- rT} (S_0 e^{rT} - K) = S_0 - K e^{- rT}$
  (using $E [S_T] = S_0 e^{rT} E [e^{X_T}] = S_0 e^{rT}$ from
  Lemma~\ref{lem:driftless}), so $\tmop{Put} = C - S_0 + K e^{- rT} = S_0
  (P_1 - 1) + K e^{- rT} (1 - P_2) = K e^{- rT} (1 - P_2) - S_0 (1 - P_1)$.
\end{proof}

\begin{corollary}
  [Pad{\'e} call/put with certified real-axis evaluation]\label{cor:callputM}
  For every $M$ and $u \in \R \setminus \{ 0 \}$, if $B_M (u, T^{\alpha}) \ne
  0$ then $\varphi_M (u, T) = \exp \left( T \hspace{0.17em} A_M (u,
  T^{\alpha}) / B_M (u, T^{\alpha}) \right)$ is defined by direct evaluation,
  with rigorous ball-arithmetic enclosures by part (iii) of
  Theorem~\ref{thm:localdata}; by Theorem~\ref{thm:realray} applied at $a =
  u$ (real), $B_M (u, T^{\alpha}) \ne 0$ for all $M \ge M_0 (u,
  T^{\alpha})$. Define $P_2^{[M]}, P_1^{[M]}, C_M, \tmop{Put}_M$ by
  substituting $\varphi_M$ for $\varphi_T$ throughout
  Theorem~\ref{thm:callput} (on the two contours $\Im  a \in \{ 0, - 1 \}$).
  If $\varphi_M \to \varphi_T$ pointwise on these contours
  (Theorem~\ref{thm:priceconv}, first part of (H1), which holds
  unconditionally at $\eta_0 = 0$ by Theorem~\ref{thm:priceconv}'s
  refinement above) and is dominated there in the sense of the analogue of
  \tmtextup{(H2)} for the kernel $K (u) = \left( 2 \pi iu \right)^{- 1} e^{-
  iuk}$, then $C_M \to C$ and $\tmop{Put}_M \to \tmop{Put}$.
\end{corollary}

\begin{proof}
  Zero-freeness of $B_M (u, \cdot)$ at $x = T^{\alpha}$ for real $u$: by
  Lemma~\ref{lem:GMdata}, $B_M (u, x) = \tilde{Q}^{(1)} (u, x)  \hspace{0.17em}
  \tilde{Q}^{(2)} (u, x)$, and each $\tilde{Q}^{(i)}$ is, up to the nonzero
  scalar $H_M^{(i)} (u)$ (Lemma~\ref{lem:clearing}), the reciprocal-flipped
  denominator $\hat{Q}^{(i)} (M, \cdot)$ of a diagonal Pad{\'e} approximant
  covered by Theorem~\ref{thm:integ}; by Theorem~\ref{thm:realray} (applied
  to the corresponding reweighted moment functional at the real parameter
  $a = u$) and Theorem~\ref{thm:stahl}, $\hat{Q}^{(i)} (M, \cdot)$ has no
  zero on $[0, T^{\alpha}]$ for $M$ at least some $M_0^{(i)} (u,
  T^{\alpha})$; take $M_0 \assign \max (M_0^{(1)}, M_0^{(2)})$. Well-definedness
  and the ball-arithmetic enclosure are then Theorem~\ref{thm:localdata}(iii)
  applied at $a = u$ with the empty pole set on $[0, T^{\alpha}]$
  (direct composition, part (i)). The convergence statement is
  Theorem~\ref{thm:callput} applied with $\varphi_M$ in place of $\varphi_T$
  in $P_1^{[M]}, P_2^{[M]}$, combined with dominated convergence on each of
  the two contours exactly as in the proof of Theorem~\ref{thm:priceconv}
  (the kernel $K (u) = (2 \pi iu)^{- 1} e^{- iuk}$, extended continuously at
  $u = 0$ as in Theorem~\ref{thm:callput}, plays the role of the abstract
  kernel of Definition~\ref{def:lewis} there).
\end{proof}

\begin{proposition}
  [Cumulants in the algebra]\label{prop:cumulants} Every $c_k (a)$, and hence
  every $g_j (a)$, is a polynomial in $a$ (Lemma~\ref{lem:polya},
  Lemma~\ref{lem:neighboring}), and for $0 \le t < \rho_0 (\epsilon)^{1 /
  \alpha}$, uniformly for $|a | \le \epsilon$ (Lemma~\ref{lem:radius}), the
  cumulants of $X_T$ are the absolutely convergent series
  \[ \kappa_n (t) \hspace{0.27em} = \hspace{0.27em} \partial_a^n L (a, t) |_{a
     = 0} \hspace{0.27em} = \hspace{0.27em} t \sum_{j \ge 0} (\partial_a^n
     g_j) (0)  \hspace{0.17em} t^{\alpha j}, \]
  with each $(\partial_a^n g_j) (0)$ a finite exact expression obtained by
  Leibniz-differentiating the M{\"u}ntz--Tau recurrence of
  Lemma~\ref{lem:muntz} in $a$. Consequently the cumulant data feeding any
  transform- or expansion-based pricing representation lies in
  $\C$-linear combinations of the same neighboring-pair objects, computed
  without approximation.
\end{proposition}

\begin{proof}
  Polynomiality of $c_k (a)$ is Lemma~\ref{lem:polya}; $g_j (a)$ is then
  polynomial in $a$ by its explicit formula in Lemma~\ref{lem:neighboring}
  (a $\C$-linear combination, with constant scalar coefficients, of $c_j (a)$
  and $c_{j + 1} (a)$). For $|a | \le \epsilon$ and $t < \rho_0
  (\epsilon)^{1 / \alpha}$, $x = t^{\alpha} < \rho_0 (\epsilon)$, so
  $\sum_j |g_j (a) |  \hspace{0.17em} x^j < \infty$ uniformly in $a$ by
  Lemma~\ref{lem:radius} together with the explicit formula for $g_j$ (each
  term of $g_j$ is a bounded scalar times $c_j$ or $c_{j + 1}$, whose sums
  $\sum |c_j (a) |  \rho^j$ are already controlled there); a locally
  uniformly convergent series of polynomials in $a$ may be differentiated
  termwise in $a$ any number of times (Weierstrass), giving the display.
  Each $\partial_a^n g_j (0)$ is, by Lemma~\ref{lem:neighboring}, the same
  $\C$-linear combination (with the same constant scalars $\theta / (j
  \alpha + 1)$, $V_0 \Gamma ((j + 1) \alpha + 1) / \Gamma (j \alpha + 2)$) of
  $\partial_a^n c_j (0)$ and $\partial_a^n c_{j + 1} (0)$; these obey the
  recurrence obtained by applying the Leibniz rule to the M{\"u}ntz--Tau
  recurrence of Lemma~\ref{lem:muntz} — $\partial_a^n c_{k + 1} = \frac{\Gamma
  (k \alpha + 1)}{\Gamma ((k + 1) \alpha + 1)} \sum_{i = 0}^n \binom{n}{i}
  \left( (\partial_a^i Q)  (\partial_a^{n - i} c_k) + (\partial_a^i R)
  \sum_{p + q = k} \sum_{l = 0}^{n - i} \binom{n - i}{l}
  (\partial_a^l c_p) (\partial_a^{n - i - l} c_q) \right)$ — closed and
  finite at every order because $P, Q, R$ are polynomials in $a$ (so all
  $a$-derivatives beyond their degree vanish).
\end{proof}

\begin{algorithm}
  [Coefficient-algebra evaluation layer for arb4j]\label{alg:arb4j} Input:
  $(A_M, B_M) \in \C [a] [x]^2$ from the Chebyshev--Wheeler pipeline; contour
  height $\eta_0$; maturity $T$; working precision; truncation bandwidth $K$.
  \begin{enumerate}
    \item For each quadrature node $a = u + \mathrm{i} \eta_0$: evaluate $B_M 
    (a, x)$ on $[0, T^{\alpha}]$ in ball arithmetic; certify zero-freeness
    (interval Descartes/argument principle). Cost $O (M)$ per node.
    
    \item Price by fixed-contour quadrature of $K (u) \exp \left( T
    \hspace{0.17em} A_M / B_M \right)$ (Theorem~\ref{thm:priceconv}); rigorous
    tails from the kernel bound. Cost: nodes $\times O (M)$.
    
    \item (Diagnostics / analytic continuation off the ray.) For each required
    pole $\xi_i$: partial fractions to get $c^{(i)}$ ($O (M^2)$ exact ops);
    $\Lambda$-sequence by \eqref{eq:universal} ($O (mK)$ ops); Taylor data
    $b^{(i)}_l$ of $e^{H_i}$ by \eqref{eq:Prec} ($O ((d_U + d_V) K)$ ops);
    assemble $A^{(i)}_n$ with certified tail via Proposition~\ref{prop:tail}.
    
    \item Cumulant channel: $\partial_a^n c_k (0)$ by the differentiated
    Gamma-ratio recurrence ($O (nM)$ per order), then $\kappa_n$ by
    Proposition~\ref{prop:cumulants}.
  \end{enumerate}
  All steps are exact recurrences or ball-arithmetic evaluations with explicit
  error terms; no Hankel matrix is ever formed or inverted (Theorem~\ref{thm:CW}),
  matching the stability profile established there.
\end{algorithm}

\section{Open problems}\label{sec:open}

\begin{problem}
  [Domination]\label{prob:H2} Prove hypothesis \tmtextup{(H2)} of
  Theorem~\ref{thm:priceconv} (equivalently, its analogue in
  Corollary~\ref{cor:callputM} for the undamped Gil--Pelaez kernel of
  \S\ref{subsec:callput}): a uniform-in-$M$ integrable majorant for $|
  \varphi_M |$ on the pricing contour, e.g.\ via a uniform-in-$a$ version of
  the Stahl rate of Theorem~\ref{thm:stahl} combined with the quadratic
  leading-order decay of $g_0 (a)$ identified in Remark~\ref{rem:H2}. What is
  {\tmem{not}} open is real-axis pole-freeness of $B_M$ itself
  (Theorem~\ref{thm:realray}, used directly in Corollary~\ref{cor:callputM});
  the remaining gap is purely the uniform-in-$M$ decay rate of $| \varphi_M
  |$.
\end{problem}

\begin{problem}
  [Resurgent structure]\label{prob:resurgence} $e^{\Pi}$ is the exponential
  factor of the Hukuhara--Levelt--Turrittin normal form of the rank-one
  connection $d - dR$ at the irregular singularity $a$
  {\tmem{({\cite[Ch.~3]{vdPS}})}}; this much is classical. Construct the
  action of the alien derivations of {\'E}calle on the $\CA_m$-module of
  Theorem~\ref{thm:module} --- the natural conjecture, suggested by
  Lemma~\ref{lem:derivation}, is that singular directions act through the
  lowering operators $\partial_{c_j}$ --- and determine the Stokes data of
  $\varphi_M$ in terms of the $c^{(i)}$ {\tmem{(framework: Mitschi--Sauzin
  {\cite{MitschiSauzin}}; Costin {\cite{Costin}})}}. No such statement is
  asserted here as a theorem.
\end{problem}

\begin{problem}
  [Exponential periods]\label{prob:motives} $e^R  \hspace{0.17em} dz$ is an
  exponential-period integrand; the category of exponential motives
  {\tmem{(Fres{\'a}n--Jossen {\cite{FresanJossen}})}} attaches cohomological
  invariants to $(\mathbb{A}^1, R)$. Position $\CA_m$ relative to the de Rham
  realization of the rank-one exponential connection: is the coefficient
  module of Theorem~\ref{thm:module} a concrete model of a known object there?
\end{problem}

\begin{problem}
  [$q$-lift]\label{prob:hall} Lift $(\Lambda_N)$ along the Hall-algebra
  isomorphism of Remark~\ref{rem:hall} to the Hall--Littlewood level: identify
  the $q$-analogue of the universal recurrence and its analytic meaning, if
  any, for $q$-deformed exponentials of principal parts.
\end{problem}

\begin{problem}
  [Non-reduction]\label{prob:nonreduction} Prove (or refute) that for generic
  $(c_1, c_2, c_3)$ the sequence $(\Lambda_N (c_1, c_2, c_3))_N$ is not
  expressible as $\left( \rho^N  \hspace{0.17em} F (\sigma N + \tau)
  \right)_N$ for any single classical hypergeometric-type function $F$ of one
  variable --- i.e. make Remark~\ref{rem:scope} a theorem, plausibly via
  differential Galois groups of the order-$3$ holonomic equation from
  \eqref{eq:universal}.
\end{problem}

\begin{thebibliography}{99}
  \bibitem{Appell}P.~Appell, J.~Kamp{\'e} de F{\'e}riet, {\tmem{Fonctions
  hyperg{\'e}om{\'e}triques et hypersph{\'e}riques. Polynomes d'Hermite}},
  Gauthier-Villars, Paris, 1926.
  
  \bibitem{Barry}P.~Barry, {\tmem{Riordan Arrays: A Primer}}, Logic Press,
  2016.
  
  \bibitem{BLL}F.~Bergeron, G.~Labelle, P.~Leroux, {\tmem{Combinatorial
  Species and Tree-like Structures}}, Encyclopedia Math. Appl. 67, Cambridge
  Univ. Press, 1998.
  
  \bibitem{CarrMadan}P.~Carr, D.~Madan, Option valuation using the fast
  Fourier transform, {\tmem{J. Comput. Finance}} \tmtextbf{2} (1999), 61--73.
  
  \bibitem{Comtet}L.~Comtet, {\tmem{Advanced Combinatorics}}, Reidel,
  Dordrecht, 1974.
  
  \bibitem{Costin}O.~Costin, {\tmem{Asymptotics and Borel Summability}},
  Chapman \& Hall/CRC, 2009.
  
  \bibitem{DLMF}NIST {\tmem{Digital Library of Mathematical Functions}},
  {\textsection}10.35, https://dlmf.nist.gov/.

  \bibitem{ElEuchRosenbaum}O.~El Euch, M.~Rosenbaum, The characteristic
  function of rough Heston models, {\tmem{Math. Finance}} \tmtextbf{29}
  (2019), 3--38.
  
  \bibitem{FGB}H.~Figueroa, J.~M.~Gracia-Bond{\'\i}a, Combinatorial Hopf
  algebras in quantum field theory I, {\tmem{Rev. Math. Phys.}} \tmtextbf{17}
  (2005), 881--976.
  
  \bibitem{FS}P.~Flajolet, R.~Sedgewick, {\tmem{Analytic Combinatorics}},
  Cambridge Univ. Press, 2009.
  
  \bibitem{FresanJossen}J.~Fres{\'a}n, P.~Jossen, {\tmem{Exponential
  Motives}}, book draft, http://javier.fresan.perso.math.cnrs.fr/.
  
  \bibitem{GilPelaez}J.~Gil-Pelaez, Note on the inversion theorem,
  {\tmem{Biometrika}} \tmtextbf{38} (1951), 481--482.
  
  \bibitem{Gragg}W.~B.~Gragg, The Pad{\'e} table and its relation to certain
  algorithms of numerical analysis, {\tmem{SIAM Review}} \tmtextbf{14} (1972),
  1--62.
  
  \bibitem{KauersPaule}M.~Kauers, P.~Paule, {\tmem{The Concrete
  Tetrahedron}}, Springer, 2011.
  
  \bibitem{KilbasSaigo}A.~A.~Kilbas, M.~Saigo, {\tmem{H-Transforms: Theory
  and Applications}}, CRC Press, 2004.
  
  \bibitem{Koutschan}C.~Koutschan, Advanced applications of the holonomic
  systems approach, PhD thesis, RISC, J.~Kepler Univ. Linz, 2009.
  
  \bibitem{Lewis2001}A.~L.~Lewis, A simple option formula for general
  jump-diffusion and other exponential L{\'e}vy processes, OptionCity working
  paper, 2001.
  
  \bibitem{Lipshitz}L.~Lipshitz, The diagonal of a D-finite power series is
  D-finite, {\tmem{J. Algebra}} \tmtextbf{113} (1988), 373--378.
  
  \bibitem{Macdonald}I.~G.~Macdonald, {\tmem{Symmetric Functions and Hall
  Polynomials}}, 2nd ed., Oxford Univ. Press, 1995.
  
  \bibitem{MilnorMoore}J.~Milnor, J.~Moore, On the structure of Hopf
  algebras, {\tmem{Ann. of Math.}} \tmtextbf{81} (1965), 211--264.
  
  \bibitem{MitschiSauzin}C.~Mitschi, D.~Sauzin, {\tmem{Divergent Series,
  Summability and Resurgence I}}, Lecture Notes in Math. 2153, Springer, 2016.
  
  \bibitem{Roman}S.~Roman, {\tmem{The Umbral Calculus}}, Academic Press,
  1984.
  
  \bibitem{Shapiro}L.~W.~Shapiro, S.~Getu, W.-J.~Woan, L.~C.~Woodson, The
  Riordan group, {\tmem{Discrete Appl. Math.}} \tmtextbf{34} (1991), 229--239.
  
  \bibitem{Stahl}H.~Stahl, The convergence of Pad{\'e} approximants to
  functions with branch points, {\tmem{J. Approx. Theory}} \tmtextbf{91}
  (1997), 139--204.
  
  \bibitem{StanleyDF}R.~P.~Stanley, Differentiably finite power series,
  {\tmem{European J. Combin.}} \tmtextbf{1} (1980), 175--188.
  
  \bibitem{vdPS}M.~van der Put, M.~F.~Singer, {\tmem{Galois Theory of Linear
  Differential Equations}}, Grundlehren Math. Wiss. 328, Springer, 2003.
  
  \bibitem{Wright1933}E.~M.~Wright, On the coefficients of power series
  having exponential singularities, {\tmem{J. London Math. Soc.}} \tmtextbf{8}
  (1933), 71--79.
\end{thebibliography}

\end{document}
